\documentclass[11pt,reqno]{amsart}
\usepackage[T1]{fontenc}
\usepackage{lmodern}
\usepackage{amsmath,amssymb,amsthm,mathtools,mathrsfs}
\usepackage[margin=1.05in]{geometry}
\usepackage{microtype,enumitem,xcolor}
\usepackage[colorlinks=true,linkcolor=blue!55!black,
 citecolor=blue!55!black,urlcolor=blue!55!black]{hyperref}
\hypersetup{pdftitle={Optimal estimates for direct images: product descent and top-degree rigidity},
 pdfauthor={Zhijian Chen, Zhuo Liu, Chenghao Qing}}
\allowdisplaybreaks[2]
\numberwithin{equation}{section}
\newtheorem{theorem}{Theorem}[section]
\newtheorem{proposition}[theorem]{Proposition}
\newtheorem{lemma}[theorem]{Lemma}
\newtheorem{corollary}[theorem]{Corollary}
\theoremstyle{definition}
\newtheorem{definition}[theorem]{Definition}
\newtheorem{example}[theorem]{Example}
\theoremstyle{remark}
\newtheorem{remark}[theorem]{Remark}
\newcommand{\OO}{\mathcal O}
\newcommand{\ddbar}{\partial\bar\partial}
\newcommand{\Id}{\operatorname{Id}}
\title[Optimal estimates and top-degree rigidity]
{Optimal estimates for direct images:\\ product descent and top-degree rigidity}
\author{Zhijian Chen}
\address{Institute of Mathematics, Academy of Mathematics and Systems Science,
 Chinese Academy of Sciences, Beijing 100190, China}
\email{chenzhijian@amss.ac.cn}
\author{Zhuo Liu}
\address{Beijing Institute of Mathematical Sciences and Applications;
 Department of Mathematics and Yau Mathematical Sciences Center,
 Tsinghua University, Beijing, China}
\email{liuzhuo@amss.ac.cn}
\author{Chenghao Qing}
\address{Yau Mathematical Sciences Center, Tsinghua University,
 Beijing 100084, China}
\email{qingchenghao@amss.ac.cn}
\date{30 October 2026}
\subjclass[2020]{Primary 32L10; Secondary 32G05, 32L20, 32W05}
\keywords{Higher direct images, optimal $L^2$ estimates, singular Hermitian
 metrics, multiplier ideals, harmonic metrics, relative duality}
\begin{document}
\begin{abstract}
We prove product descent of optimal $L^2$ estimates for prescribed
direct-image metrics.  For arbitrary singular semipositive line metrics,
a rigidity theorem gives non-product adjoint descent in top fiber degree $q=n$.
We identify the smooth adjoint metric with the Mourougane--Takayama
metric and give counterexamples to unrestricted descent.
\end{abstract}
\maketitle
\enlargethispage{8pt}
\tableofcontents

\section{Introduction}\label{sec:introduction}

Let $\pi:X\to Y$ be a proper holomorphic submersion between complex
manifolds, with compact
fibers $X_y$ of dimension $n$ and base dimension $m$.  Fix K\"ahler
metrics and a Hermitian holomorphic line bundle $(L,h)$ on $X$.
For smooth $h$, the fiber cohomology
$H^q(X_y,\Omega_{X_y}^p\otimes L_y)$ carries its harmonic norm.
For a metric with plurisubharmonic weights, the relevant sheaf is
\[
 \mathcal F^{p,q}_h
 =R^q\pi_*(\Omega_{X/Y}^p\otimes L\otimes\mathcal I(h)),
\]
and the prescribed norm is the minimum weighted $L^2$ norm of a
Dolbeault representative, wherever the fiberwise identification is
defined.  The multiplier ideal and the choice of norm are part of the
statement.

The positivity of direct images developed from Hodge-theoretic
semipositivity, the theorems of Fujita and Kawamata
\cite{Fuj78,Kaw81}, Viehweg's weak positivity \cite{Vie83}, and
Koll\'ar's higher-direct-image theorems \cite{Kol86a,Kol86b}.
Takegoshi treated Nakano semipositive coefficients under proper
K\"ahler morphisms using harmonic methods \cite{Take95}.
Mourougane--Takayama studied Hodge metrics on adjoint direct images
\cite{MT07}; Berndtsson proved Nakano semipositivity of the natural
$L^2$ metric on $f_*(K_{X/Y}\otimes L)$ for a smooth semipositive
line metric \cite{Ber09}.
Mourougane--Takayama obtained the higher-direct-image version
\cite{MT08} and extended the induced tautological metrics across
singular fibers \cite{MT09}.
Curvature formulas for other Hodge degrees were subsequently developed
by, among others, Liu--Yang, Naumann, Berndtsson--P\u aun--Wang,
and Choi--Schumacher \cite{LY14,Nau21,BPW22,CS24}.

For singular coefficients, Berndtsson--P\u aun established positivity
of relative Bergman metrics \cite{BP08}, and P\u aun--Takayama
developed singular metric positivity for twisted pluricanonical
direct images \cite{PT18}.
These methods enter Cao--P\u aun's work on algebraic fiber spaces over
abelian varieties \cite{CP17}; Hacon--Popa--Schnell give a detailed
treatment of the adjoint direct-image metric and its minimal extension
property \cite{HPS18}.
The formulation of positivity by optimal $\bar\partial$ estimates
was developed in \cite{DNWZ23,WatCurv}.
For ordinary adjoint direct images, singular Nakano positivity is proved
in \cite{Wat25,IMW24}; the latter allows arbitrary singular
semipositive line metrics and proper K\"ahler fibrations.
Takayama proves singular Griffiths semipositivity and the minimal
extension property for higher direct images with smooth Nakano
semipositive coefficients \cite{Tak24}.

The multiplier ideal in a singular higher direct image is essential.
Demailly--Peternell--Schneider prove a hard Lefschetz theorem with
multiplier ideals \cite{DPS01}, and Wu studies the resulting
parallel sections \cite{Wu21}.
Matsumura establishes relative injectivity and torsion-freeness for
arbitrary singular semipositive metrics \cite{Mat22}.
Our top-degree argument uses instead the rigidity forced by
nonvanishing, together with relative duality and the divisorial
decomposition of a positive current \cite{RRV71,Siu74}.

The product descent question compares an optimal estimate in total
degree $(a+p,b+q)$ with the estimate in base degree $(a,b)$ for the
prescribed direct-image metric.  The smooth and weighted optimal
estimates are defined in Definitions~\ref{def:smooth-optimal}
and~\ref{def:singular-estimate}; $\mathscr E(H)$ denotes the sheaf of
locally square-integrable holomorphic sections.  We use the ranges
\begin{equation}\label{eq:ranges}
\begin{array}{ll}
\mathrm{(i)}&a=m,\quad p=n,\quad 1\le b\le m,\quad0\le q\le n,\\
\mathrm{(ii)}&b=m,\quad q=n,\quad 1\le a\le m,\quad0\le p\le n.
\end{array}
\end{equation}

\begin{theorem}[Smooth product descent]\label{thm:product-descent}
Let $(Y,\omega_Y)$ be complete K\"ahler of dimension $m\ge1$, admitting
a positive Hermitian holomorphic line bundle.  Let $(Z,\omega_Z)$
be smooth projective of dimension $n$, and set
\[
 X=Y\times Z,\quad\pi=\operatorname{pr}_Y,\quad
 \omega_X=\omega_Y+\omega_Z,\quad
 (L,h)=(G,h_G)\boxtimes(M,h_M),
\]
where $G$ and $M$ are holomorphic line bundles on $Y$ and $Z$,
respectively, and both metrics are smooth.  In either range \eqref{eq:ranges}, if $(L,h)$
satisfies the optimal $(a+p,b+q)$ estimate, then
\[
 \bigl(\mathcal F,H\bigr)
 =\bigl(R^q\pi_*(\Omega_{X/Y}^p\otimes L),H\bigr)
\]
satisfies the optimal $(a,b)$ estimate for its prescribed harmonic
metric.
\end{theorem}

In range~\textnormal{(ii)}, Proposition~\ref{thm:variable-metric}
allows an arbitrary smooth metric on $G\boxtimes M$, with the same
product geometry.  Singular semipositive product metrics give
the following result.

\begin{theorem}[Singular product descent]\label{thm:singular-product}
Let $(Y,\omega_Y)$ be complete K\"ahler of dimension $m\ge1$
and $(Z,\omega_Z)$ compact K\"ahler of dimension $n$.  On $X=Y\times Z$, with $\pi=\operatorname{pr}_Y$, let
\[
 (L,h)=(G,h_G)\boxtimes(M,h_M),
\]
where $G$ and $M$ are holomorphic line bundles on $Y$ and $Z$
with arbitrary psh metric weights.  Fix $0\le p,q\le n$ and define
\[
 V=H^q(Z,\Omega_Z^p\otimes M\otimes\mathcal I(h_M)),\quad
 E_0=G\otimes V,\quad H=h_G\otimes H_V.
\]
Here $H_V$ is the minimum weighted $L^2$ norm of Dolbeault
representatives, whose existence is proved in
Lemma~\ref{prop:fiber-norm}.  There is a canonical identification
\begin{equation}\label{eq:singular-product-sheaf}
 \begin{aligned}
 \mathcal F&:=R^q\pi_*(\Omega_{X/Y}^p\otimes L\otimes\mathcal I(h))\\
  &=(G\otimes\mathcal I(h_G))\otimes V=\mathscr E(H).
 \end{aligned}
\end{equation}
For almost every $y$, $H_y$ is the prescribed weighted fiberwise
quotient norm.  The pair in \eqref{eq:singular-product-sheaf}
satisfies the weighted optimal estimates in all degrees
\[
 (m,b),\quad1\le b\le m,\qquad\text{and}\qquad
 (a,m),\quad1\le a\le m,
\]
for every $p,q$.  The metric $H$ on the torsion-free sheaf
$\mathcal F$ is Griffiths semipositive and has the minimal extension
property.
\end{theorem}

The top adjoint degree admits a non-product result for completely
arbitrary singular semipositive coefficient metrics.

\begin{theorem}[Rigidity and descent in top adjoint degree]
\label{thm:singular-top-rigidity}
Let $\pi:X\to Y$ be a proper surjective holomorphic submersion
between connected complex manifolds, with connected fibers of dimension
$n\ge1$.  Fix a K\"ahler form $\omega_X$ on $X$, put
$\omega_y=\omega_X|_{X_y}$, and let $(L,h)$ be a holomorphic line
bundle with arbitrary psh metric weights.  Set
\[
 \mathcal F_n=R^n\pi_*
   (K_{X/Y}\otimes L\otimes\mathcal I(h)).
\]
Either $\mathcal F_n=0$, or there exist an effective
divisor $D$ on $X$ and a holomorphic line bundle $G$ on $Y$ with a singular
Hermitian metric $h_G$ having psh local weights such that
\begin{equation}\label{eq:top-singular-factorization}
 (L,h)\simeq
 (\mathcal O_X(D),h_D)\otimes\pi^*(G,h_G).
\end{equation}
Here $s_D$ is the canonical section of $\mathcal O_X(D)$ and $h_D$
is the divisorial metric characterized by $|s_D|_{h_D}=1$ off $D$.
In the second case,
\begin{equation}\label{eq:top-singular-sheaf}
 \mathcal I(h)=\mathcal O_X(-D)\cdot
       \pi^{-1}\mathcal I(h_G)\mathcal O_X,
 \qquad
 \mathcal F_n\simeq G\otimes\mathcal I(h_G).
\end{equation}
The prescribed weighted fiberwise quotient norm is, almost everywhere,
\begin{equation}\label{eq:top-singular-metric}
 H=V^{-1}h_G,
 \qquad
 V=\int_{X_y}\frac{\omega_y^n}{n!}>0,
\end{equation}
where $V$ is independent of $y$, and the trace of a top Dolbeault class
is normalized by integration. In particular,
$\mathcal F_n=\mathscr E(H)$ as a subsheaf of $G$.

If $(Y,\omega_Y)$ is complete K\"ahler of dimension $m\geq1$, then
$(\mathcal F_n,H)$ satisfies the singular optimal $(m,b)$ estimate for
every $1\leq b\leq m$.
\end{theorem}

For smooth vector coefficients, isometric Serre duality gives the
following complementary statement.

\begin{theorem}[Top direct images under infinitesimal triviality]
\label{thm:top-dual}
Let $\pi:X\to Y$ be a proper surjective holomorphic submersion of
relative dimension $n$ between K\"ahler manifolds.  Assume that all
fiberwise Kodaira--Spencer maps vanish.  Let $(E,h)$ be a smooth
Hermitian holomorphic vector bundle such that $(E^*,h^*)$ is Nakano
seminegative, and assume that $h^n(X_y,E_y)$ is locally constant.
Then the prescribed harmonic metric on $R^n\pi_*E$ is smooth and
dual Nakano semipositive, that is, its dual metric is Nakano
seminegative.
\end{theorem}

For smooth semipositive line coefficients, Mourougane--Takayama
\cite[Theorem~1.1]{MT08} imply that the prescribed harmonic metric on
\[
 R^q\pi_*(K_{X/Y}\otimes L),\qquad 0\le q\le n,
\]
is Nakano semipositive for every proper K\"ahler submersion.
On a complete K\"ahler base this gives the optimal $(m,b)$ estimates
for $1\le b\le m$.
Remark~\ref{rem:mt-comparison} gives the detailed metric comparison.

The smooth split-product implication is a consequence of the curvature
characterization in \cite{WatCurv} and the fiberwise Bochner identity.
The singular split-product result combines scalar weighted estimates
with the multiplier-ideal product formula \cite{Dem82,DEL00}.
We therefore present these statements as consequences of existing
theory.
The failure of partial positivity under fiber integration has a close
precedent in Inayama's polydisc counterexamples
\cite[Theorem~1.1]{Ina26}; our compact-fiber examples concern a proper
direct image with its prescribed cohomological metric.
The top-degree singular result concerns $q=n$ and identifies the
actual sheaf and its quotient metric.

Section~\ref{sec:counterexamples} gives compact-fiber counterexamples
to the unrestricted degreewise implication.  One occurs on a trivial
family in range~\textnormal{(i)}, with a nonsplit metric which is not
semipositive; another uses a flat coefficient on a varying elliptic
family in range~\textnormal{(ii)}.  Thus infinitesimal triviality alone
does not replace metric splitting in the original degreewise problem.
Section~\ref{sec:extension} proves Griffiths semipositivity and the
minimal extension property for the scalar-factor metrics in the
singular positive results.


\section{Curvature estimates and weighted cohomology}\label{sec:analytic}

We use $dV_\omega=\omega^d/d!$ on a K\"ahler $d$-fold and the
curvature convention $i\Theta_{L,e^{-\varphi}}=i\partial\bar\partial\varphi$.
All estimates use the specified form and coefficient metrics.
A singular semipositive line metric
means $h=e^{-\varphi}$ in local holomorphic frames, with
$\varphi\in\operatorname{PSH}$ and $\varphi\not\equiv-\infty$.
It is positive and finite almost everywhere and locally bounded below
by a positive smooth metric.  Put
\[
 C_{E;r,s}=[i\Theta_{E,h}\otimes\Id,\Lambda_\omega]_{r,s}.
\]
For scalar curvature with eigenvalues $\lambda_1,\ldots,\lambda_d$,
its eigenvalue on $dz^I\wedge d\bar z^J$ is
\begin{equation}\label{eq:commutator-eigenvalue}
 \sum_{i\in I}\lambda_i+\sum_{j\in J}\lambda_j
       -\sum_{k=1}^d\lambda_k.
\end{equation}

\begin{definition}\label{def:smooth-optimal}
Let $(S,\omega)$ admit a positive Hermitian holomorphic line bundle.
In a degree $r=d$, $1\le s\le d$, or $s=d$, $1\le r\le d$,
a smooth Hermitian bundle $(E,h)$ satisfies the \emph{optimal $(r,s)$
estimate} if, for every smooth positive auxiliary line bundle $(A,h_A)$
and every $\bar\partial$-closed $f\in C_c^\infty(S,\Lambda^{r,s}T_S^*
\otimes E\otimes A)$, there is a distributional solution
$u\in L^2_{r,s-1}(S,E\otimes A)$ with
\begin{equation}\label{eq:optimal-bound}
 \bar\partial u=f,\qquad
 \|u\|_{h h_A,\omega}^2\le
 \int_S\langle B_A^{-1}f,f\rangle_{h h_A,\omega}\,dV_\omega,
 \quad B_A=[i\Theta_A\otimes\Id_E,\Lambda_\omega]_{r,s}.
\end{equation}
\end{definition}
Formula~\eqref{eq:commutator-eigenvalue} gives $B_A>0$ in these ranges.

\begin{proposition}[Smooth curvature characterization]\label{prop:smooth-test}
On a complete K\"ahler manifold admitting a positive Hermitian line
bundle, Definition~\ref{def:smooth-optimal} is equivalent to
$C_{E;r,s}\ge0$.
\end{proposition}
The necessary implication is \cite[Theorems 1.5--1.6]{WatCurv}, with the
fixed-metric definitions in \cite[Definitions 1.3--1.4]{WatCurv}.
The sufficient implication follows from the next lemma.

For a nonnegative Hermitian operator $C$, define its extended inverse
quadratic form by
\begin{equation}\label{eq:extended-inverse}
 q_C(f)=\sup_v\{2\operatorname{Re}\langle f,v\rangle
                         -\langle Cv,v\rangle\}.
\end{equation}
To see its meaning, diagonalize $C$ at a point, with eigenvalues
$\mu_\nu\ge0$, and write $f=\sum_\nu f_\nu e_\nu$.  Maximizing each
scalar expression in \eqref{eq:extended-inverse} gives
\[
 q_C(f)=
 \begin{cases}
 \displaystyle\sum_{\mu_\nu>0}\frac{|f_\nu|^2}{\mu_\nu},
       &f_\nu=0\text{ whenever }\mu_\nu=0,\\[4pt]
 +\infty,&\text{otherwise}.
 \end{cases}
\]
Thus it equals $\langle C^\dagger f,f\rangle$ for
$f\perp\ker C$, where $C^\dagger$ is the inverse on $(\ker C)^\perp$
and zero on $\ker C$.  This definition also works when the rank of
$C$ changes.  Directly comparing the expressions inside the supremum gives
\begin{equation}\label{eq:inverse-order}
 C\ge B\ge0\quad\Longrightarrow\quad q_C\le q_B.
\end{equation}

\begin{lemma}[Complete K\"ahler estimate]\label{lem:complete-estimate}
Let $(S,\omega)$ be complete K\"ahler, let $(E,h)$ be smooth, and let
$C=C_{E;r,s}\ge0$, where $s\ge1$.  If $f\in L^2_{r,s}(S,E)$,
$\bar\partial f=0$, and $I(f)=\int_Sq_C(f)\,dV_\omega<\infty$, then
there exists $u\in L^2_{r,s-1}(S,E)$ satisfying
$\bar\partial u=f$ and $\|u\|^2\le I(f)$.
\end{lemma}
\begin{proof}
All operators in this proof act between the Hilbert spaces defined by
the fixed metrics.  Set
\[
 T=\bar\partial_{\max}:L^2_{r,s-1}(S,E)\longrightarrow L^2_{r,s}(S,E),
 \qquad
 S_1=\bar\partial_{\max}:L^2_{r,s}(S,E)\longrightarrow L^2_{r,s+1}(S,E),
\]
with $S_1=0$ if $s=d$.  They are densely defined and closed, since
convergence in their graph norms implies convergence of the
corresponding distributional derivatives.  Moreover,
$\operatorname{im}T\subset\ker S_1$ follows from
$\bar\partial^2=0$ in distributions.  The operator $T^*$ is the
Hilbert-space adjoint, whose local expression is the formal
$\bar\partial$-adjoint.

The Bochner--Kodaira--Nakano identity gives, for compactly supported
smooth $w$,
\[
 \|\bar\partial w\|^2+\|\bar\partial^*w\|^2
 =\|D'w\|^2+\|(D')^*w\|^2+
   \int_S\langle Cw,w\rangle\,dV_\omega.
\]
We first justify using its inequality for
$w\in\operatorname{Dom}S_1\cap\operatorname{Dom}T^*$.
Completeness provides compactly supported Lipschitz cutoffs
$0\le\chi_\nu\le1$ with
\[
 \chi_\nu\longrightarrow1\quad\text{pointwise},\qquad
 |d\chi_\nu|_\omega\le C/\nu.
\]
For example, compose the distance from a fixed point with a cutoff
that equals one on $[0,\nu]$ and zero on $[2\nu,\infty)$.
On a disconnected manifold one first truncates to finitely many
components and then makes the same construction on each component.
The distributional product formulas are
\[
 S_1(\chi_\nu w)=\chi_\nu S_1w+\bar\partial\chi_\nu\wedge w,
 \qquad
 T^*(\chi_\nu w)=\chi_\nu T^*w
                    -(\bar\partial\chi_\nu\wedge)^*w.
\]
The error terms have norms at most $C\nu^{-1}\|w\|$.  Dominated
convergence therefore shows that $\chi_\nu w\to w$ in the joint
graph norm
$\|w\|+\|S_1w\|+\|T^*w\|$.  For a fixed $\nu$, choose finitely many
coordinate and bundle charts covering its support.  After a smooth
partition of unity, convolution of the coefficient distributions
approximates $\chi_\nu w$ by compactly supported smooth forms in
that same graph norm.  This is the local Friedrichs lemma applied
simultaneously to the first-order operator $(S_1,T^*)$; its
coefficients and the reference metrics are smooth on the compact
coordinate sets.  A diagonal choice gives smooth $w_\nu$ such that
\[
 w_\nu\to w,\qquad S_1w_\nu\to S_1w,\qquad T^*w_\nu\to T^*w
 \quad\text{in }L^2.
\]
Choose a subsequence converging pointwise almost everywhere.
Since $C\ge0$, Fatou's lemma, followed by the smooth Bochner
inequality, gives
\[
 \int_S\langle Cw,w\rangle\,dV_\omega
 \le\liminf_\nu\int_S\langle Cw_\nu,w_\nu\rangle\,dV_\omega
 \le\|S_1w\|^2+\|T^*w\|^2.
\]
Finite $I(f)$ implies $f\perp\ker C$ almost everywhere.  The
diagonal expression for $q_C$ gives the pointwise inequality
\[
 |\langle f,w\rangle|
 \le q_C(f)^{1/2}\langle Cw,w\rangle^{1/2}.
\]
After integration and Cauchy--Schwarz, for
$w\in\ker S_1\cap\operatorname{Dom}T^*$ we obtain
\[
 |(f,w)|^2\le I(f)\int_S\langle Cw,w\rangle\,dV_\omega
             \le I(f)\|T^*w\|^2.
\]
Consequently the rule
\[
 \ell(T^*w)=(f,w),\qquad
 w\in\ker S_1\cap\operatorname{Dom}T^*,
\]
defines a bounded functional on this subspace of $L^2_{r,s-1}$:
if $T^*w=0$, the preceding inequality forces $(f,w)=0$.
Its norm is at most $I(f)^{1/2}$.  With inner products linear in
the first variable, $\ell$ is antilinear; Hahn--Banach and Riesz
give $u\in L^2_{r,s-1}$ with
$\ell(v)=(u,v)$ and $\|u\|^2\le I(f)$.

It remains to test against all of $\operatorname{Dom}T^*$.
For arbitrary
$w\in\operatorname{Dom}T^*$, write $w=w_0+w_\perp$ according to
$\ker S_1\oplus(\ker S_1)^\perp$.  Since $\operatorname{im}T\subset
\ker S_1$, for every $v\in\operatorname{Dom}T$ one has
$(Tv,w_\perp)=0$.  The definition of the Hilbert adjoint implies
$w_\perp\in\operatorname{Dom}T^*$ and $T^*w_\perp=0$.
Thus $w_0\in\operatorname{Dom}T^*\cap\ker S_1$.  Since also
$f\in\ker S_1$, we conclude
\[
 (u,T^*w)=(u,T^*w_0)=(f,w_0)=(f,w).
\]
This identity means that $u\in\operatorname{Dom}(T^*)^*$ and
$(T^*)^*u=f$.  Closedness and dense definition of $T$ give
$(T^*)^*=T$, so $\bar\partial u=f$.  See also
\cite[Chapters 4--5]{Dem12}.
\end{proof}

For singular metrics we formulate optimal estimates for finite-energy data.
\begin{definition}\label{def:singular-estimate}
Let $E_0$ be a holomorphic vector bundle on a smooth K\"ahler manifold
and $H$ a measurable metric, locally bounded below by a positive
smooth metric.  Define
\[
 \mathscr E(H)(U)=\{s\in H^0(U,E_0):|s|_H^2\in L^1_{\rm loc}(U)\}.
\]
The pair $(\mathscr E(H),H)$ satisfies the \emph{weighted optimal
$(a,b)$ estimate} if, for every smooth positive auxiliary line $A$
and every distributionally closed measurable $(a,b)$-form $f$ with
\[
 I_A(f)=\int_Y\langle B_A^{-1}f,f\rangle_{Hh_A,\omega_Y}\,dV_Y<\infty,
\]
there is a distributional solution $\bar\partial u=f$ with
$\|u\|_{Hh_A,\omega_Y}^2\le I_A(f)$.  The Dolbeault operators are
their maximal distributional realizations.
\end{definition}
Finite inverse energy implies local weighted $L^2$ integrability,
since $B_A^{-1}$ has a positive lower bound on every compact set.

\begin{lemma}[$L^2$ removability]\label{lem:l2-removal}
Let $A$ be a nowhere-dense analytic subset of a complex manifold.  If
$u,g\in L^2_{\mathrm{loc}}$ and $\bar\partial u=g$ off $A$, then
this identity holds distributionally everywhere.
\end{lemma}
\begin{proof}
Work in a coordinate neighborhood $U\subset\mathbb C^N$, with a
holomorphic bundle frame if necessary.  All calculations are
coefficientwise, and the smooth reference metrics are equivalent
on compact subsets to Euclidean metrics.  Let $u$ have type $(p,q)$
and put $k=p+q$.  If $q=N$, the equation is automatic, so assume
$q<N$ and fix $\tau\in C_c^\infty(U,\Lambda^{N-p,N-q-1})$.
Choose a nonzero holomorphic function $f$ on $U$ with
$A\cap U\subset\{f=0\}$ and $|f|<1$.  Set
\[
 \chi_j=\max\{0,1+j^{-1}\log|f|\}.
\]
These functions are locally Lipschitz and psh, and satisfy
\[
 0\le\chi_j\le1,\qquad
 \chi_j=0\text{ on }\{|f|\le e^{-j}\},\qquad
 \chi_j\longrightarrow1\text{ almost everywhere}.
\]
For $\eta\in C_c^\infty(U)$, nonnegative smooth regularizations of
$\chi_j$ and the identity
$\Delta(\chi_j^2)=2|\nabla\chi_j|^2+2\chi_j\Delta\chi_j$
give, in the distributional sense,
\[
 2\int_U\eta^2|\nabla\chi_j|^2
 \le\langle\Delta(\chi_j^2),\eta^2\rangle
 =\int_U\chi_j^2\Delta(\eta^2)
 \le\int_U|\Delta(\eta^2)|.
\]
Consequently $\sup_j\|d\chi_j\|_{L^2(K)}<\infty$ for every
$K\Subset U$.  Also $\chi_j\to1$ in $L^2_{\mathrm{loc}}$.  For
every compactly supported smooth real one-form $\zeta$,
\[
 (d\chi_j,\zeta)_{L^2}
   =(\chi_j,d^*\zeta)_{L^2}
   \longrightarrow(1,d^*\zeta)_{L^2}=0.
\]
The local $L^2$ bound and density of smooth tests imply
$d\chi_j\rightharpoonup0$ in $L^2_{\mathrm{loc}}$.

For fixed $j$, the support of $\chi_j\tau$ is compact and contained
in $\{|f|\ge e^{-j}\}\subset U\setminus A$.  Coefficientwise
mollification gives forms
$\tau_{j,\varepsilon}\in C_c^\infty(U\setminus\{f=0\})$, with
supports in a fixed compact subset of $U\setminus\{f=0\}$, such that
\[
 \|\tau_{j,\varepsilon}-\chi_j\tau\|_{W^{1,2}}\longrightarrow0
 \qquad(\varepsilon\to0).
\]
The equation off $A$ gives
$\int g\wedge\tau_{j,\varepsilon}
 =(-1)^{k+1}\int u\wedge\bar\partial\tau_{j,\varepsilon}$.
For a compact set containing these supports, Cauchy--Schwarz yields
\[
 \begin{split}
 &\left|\int g\wedge(\tau_{j,\varepsilon}-\chi_j\tau)\right|
 +\left|\int u\wedge
       \bar\partial(\tau_{j,\varepsilon}-\chi_j\tau)\right|\\
 &\hspace{1cm}\le
 C(\|g\|_{L^2}+\|u\|_{L^2})
   \|\tau_{j,\varepsilon}-\chi_j\tau\|_{W^{1,2}}
 \longrightarrow0 .
 \end{split}
\]
Letting $\varepsilon\to0$ therefore gives the exact identity
\[
 \int_U\chi_jg\wedge\tau
 =(-1)^{k+1}\left(
    \int_Uu\wedge\bar\partial\chi_j\wedge\tau
      +\int_U\chi_ju\wedge\bar\partial\tau\right).
\]
For $K\Subset U$ containing $\operatorname{supp}\tau$, the functional
\[
 \ell_\tau(\alpha)=\int_Uu\wedge\alpha\wedge\tau,\qquad
 |\ell_\tau(\alpha)|
       \le C_\tau\|u\|_{L^2(K)}\|\alpha\|_{L^2(K)}
\]
is continuous on $L^2_{0,1}(K)$.  Hence
$\ell_\tau(\bar\partial\chi_j)\to0$ by weak convergence.  Moreover,
\[
 \|(\chi_j-1)u\|_{L^2(K)}^2
   =\int_K|\chi_j-1|^2|u|^2\longrightarrow0,\qquad
 \|(\chi_j-1)g\|_{L^2(K)}\longrightarrow0,
\]
because $|\chi_j-1|\le1$ and $u,g\in L^2(K)$.
Passing to the limit in the preceding identity gives
\[
 \int_Ug\wedge\tau=(-1)^{k+1}\int_Uu\wedge\bar\partial\tau
 \qquad\text{for every }\tau,
\]
which is $\bar\partial u=g$ on $U$.
\end{proof}

\begin{lemma}[Weighted fiber cohomology]\label{prop:fiber-norm}
Let $Z$ be compact Hermitian, and let $(M,h_M)$ have psh local weights.
For every $p,q$, the vector space
\[
 V=H^q(Z,\Omega_Z^p\otimes M\otimes\mathcal I(h_M))
\]
has the positive-definite quotient norm
\begin{equation}\label{eq:weighted-quotient}
 \|v\|_{H_V}^2
 =\min\{\|\alpha\|_{h_M}^2:\bar\partial\alpha=0, [\alpha]=v\}.
\end{equation}
The minimum is attained uniquely in
$\ker\bar\partial\cap\ker\bar\partial^*$ for the maximal weighted
Hilbert complex.
\end{lemma}
\begin{proof}
Let $\mathscr L^{p,j}$ consist of germs $u$ such that both $u$ and
$\bar\partial u$ are locally $L^2(h_M)$.  The complex
\begin{equation}\label{eq:weighted-resolution}
 0\longrightarrow\Omega_Z^p\otimes M\otimes\mathcal I(h_M)
 \longrightarrow\mathscr L^{p,0}\xrightarrow{\bar\partial}
 \cdots\xrightarrow{\bar\partial}\mathscr L^{p,n}\longrightarrow0
\end{equation}
is a fine resolution, as follows.  For a smooth partition function
$\rho$, the identity
\[
 \bar\partial(\rho u)=\rho\bar\partial u+\bar\partial\rho\wedge u
\]
shows that multiplication by $\rho$ preserves each graph sheaf:
on a relatively compact coordinate set,
\[
 \|\rho u\|_{h_M}+\|\bar\partial(\rho u)\|_{h_M}
 \le C_\rho\bigl(\|u\|_{h_M}+\|\bar\partial u\|_{h_M}\bigr).
\]
Thus smooth partitions of unity act on the sheaves
$\mathscr L^{p,j}$.

For local exactness, trivialize $M$ on a coordinate ball and write
$h_M=e^{-\varphi}$.  Let $j\ge1$ and let
$u\in\mathscr L^{p,j}$ satisfy $\bar\partial u=0$.
After shrinking to a Euclidean ball $B$ with
$\|u\|_{L^2(B,e^{-\varphi})}<\infty$, write
\[
 u=\sum_{|I|=p}dz^I\wedge u_I,\qquad
 \bar\partial u_I=0.
\]
For $c>0$, the scalar H\"ormander theorem on $B$ with weight
$\varphi+c|z|^2$ gives $(0,j-1)$-forms $v_I$ satisfying
\[
 \bar\partial v_I=u_I,\qquad
 \int_B|v_I|^2e^{-\varphi-c|z|^2}
 \le\frac1{cj}\int_B|u_I|^2e^{-\varphi-c|z|^2}.
\]
The inverse bound follows from
$i\partial\bar\partial(\varphi+c|z|^2)\ge
c\,i\sum dz_\ell\wedge d\bar z_\ell$: multiply a scalar
$(0,j)$-form by the Euclidean holomorphic volume form to obtain
an $(n,j)$-form, whose curvature operator has lower bound $cj$.
Since $e^{-c|z|^2}$ is bounded above and below by positive constants
on $B$, the form
\[
 v=(-1)^p\sum_{|I|=p}dz^I\wedge v_I
\]
is locally $L^2(h_M)$ and satisfies $\bar\partial v=u$.
Therefore $v\in\mathscr L^{p,j-1}$.
The Euclidean and original smooth form metrics are locally
equivalent, so this proves exactness for the prescribed metric.
In degree zero, a distributional $\bar\partial$-closed coefficient
is holomorphic, and its weighted square integrability is exactly
membership in $\mathcal I(\varphi)$.  This identifies the kernel
with $\Omega_Z^p\otimes M\otimes\mathcal I(h_M)$.
See also \cite[Theorem 5.3(a)]{WatNadel}.

Write $\mathcal H_j=L^2_{p,j}(Z,M,h_M)$ and
$d_j=\bar\partial_{\max}$, so
\[
 \operatorname{Dom}d_j
  =\{u\in\mathcal H_j:\bar\partial u
           \text{ is represented by an element of }\mathcal H_{j+1}\}.
\]
These operators are closed.  Indeed, if
$u_\ell\to u$ in $\mathcal H_j$ and
$d_ju_\ell\to v$ in $\mathcal H_{j+1}$, the local lower bound for
$h_M$ gives the same convergences in ordinary $L^2_{\rm loc}$.
Thus $\bar\partial u=v$ in distributions, so
$u\in\operatorname{Dom}d_j$ and $d_ju=v$.

To verify dense definition, put
$A=\operatorname{Supp}(\mathcal O_Z/\mathcal I(h_M))$.
This is a nowhere-dense analytic subset.  On $Z\setminus A$, the weight
$e^{-\varphi}$ is locally integrable because
$1\in\mathcal I(\varphi)$ there.  Hence every compactly supported
smooth form in $Z\setminus A$, as well as its derivative, is
weighted $L^2$.  Smooth forms are dense for this locally finite
weighted measure on relatively compact coordinate sets.
Moreover, if $\chi_\ell$ are bounded cutoffs exhausting
$Z\setminus A$, then for every $u\in\mathcal H_j$,
\[
 \|(1-\chi_\ell)u\|_{h_M}^2
   =\int_Z|1-\chi_\ell|^2|u|_{h_M}^2\,dV_Z\longrightarrow0,
\]
since $A$ has measure zero.  A partition of unity and local
approximation prove density in all of $\mathcal H_j$.

Global sections of the graph sheaves on compact $Z$ are exactly
these maximal domains.  The fine resolution therefore computes
the algebraic, initially unreduced, Hilbert-complex quotient:
\[
 V=\ker d_q/\operatorname{im}d_{q-1},\qquad \dim V<\infty.
\]
Finiteness follows from coherence of the multiplier ideal and
finite-dimensional coherent cohomology on a compact complex
manifold.  We now prove that this algebraic quotient is separated.
For $q\ge1$, equip
$D=\operatorname{Dom}d_{q-1}$ with the graph norm
\[
 \|v\|_D^2=\|v\|_{\mathcal H_{q-1}}^2
                         +\|d_{q-1}v\|_{\mathcal H_q}^2,
 \qquad K=\ker d_q.
\]
Both $D$ and $K$ are Banach spaces, and
$d_{q-1}:D\to K$ is bounded.  Choose a finite-dimensional algebraic
complement $W\subset K$ to its range, with the norm induced from
$K$.  The bounded surjection
\[
 D\oplus W\longrightarrow K,\qquad (u,w)\longmapsto d_{q-1}u+w
\]
has kernel $\ker d_{q-1}\oplus\{0\}$.  The open mapping theorem gives
a topological isomorphism
\[
 (D/\ker d_{q-1})\oplus W\xrightarrow{\ \sim\ }K,\qquad
 ([v],w)\longmapsto d_{q-1}v+w.
\]
The first summand is closed in the domain, and its image is
$B_q=\operatorname{im}d_{q-1}$.  Therefore $B_q$ is closed in $K$.
Since $K$ is closed in $\mathcal H_q$, $B_q$ is also closed in
$\mathcal H_q$.

Let $\alpha\in K$ represent a class, and let $P_{B_q}$ be the
orthogonal projection in $\mathcal H_q$.  Then
\[
 \alpha_h=\alpha-P_{B_q}\alpha\in K\cap B_q^\perp,\qquad
 [\alpha_h]=[\alpha].
\]
The definition of the Hilbert adjoint gives
\[
 B_q^\perp=\ker d_{q-1}^*,\qquad
 \|\alpha_h+d_{q-1}v\|^2
       =\|\alpha_h\|^2+\|d_{q-1}v\|^2
       \quad(v\in\operatorname{Dom}d_{q-1}).
\]
This proves existence and uniqueness of the minimum in
\eqref{eq:weighted-quotient}.  A class has minimum zero only when
$\alpha_h=0$, which is equivalent to its being zero.  For $q=0$,
the same argument uses $B_0=\operatorname{im}d_{-1}=0$.
\end{proof}


\section{Smooth product descent}\label{sec:smooth-product}


\begin{proof}[Proof of Theorem~\ref{thm:product-descent}]
Put $V=H^q(Z,\Omega_Z^p\otimes M)$.  Product base change and the
projection formula give the holomorphic identification
$\mathcal F=G\otimes_{\mathbb C}V$.  Choose a basis
$v_1,\ldots,v_N$ of $V$ and its harmonic representatives
$\alpha_1,\ldots,\alpha_N$ for the fixed metrics $(h_M,\omega_Z)$.
If $e_G$ is a local frame of $G$, the representative of
$e_G(y)\otimes v_\nu$ on $X_y$ is $e_G(y)\otimes\alpha_\nu$:
the factor $|e_G(y)|_{h_G}^2$ is a positive constant on that fiber,
so multiplication by it preserves the minimum-norm representative.
Thus the Gram matrix is
\[
 H_{\nu\bar\mu}(y)
  =|e_G(y)|_{h_G}^2
       \int_Z\langle\alpha_\nu,\alpha_\mu\rangle_{h_M,\omega_Z}\,dV_Z
  =|e_G(y)|_{h_G}^2(H_V)_{\nu\bar\mu}.
\]
This proves the holomorphic isometry
\begin{equation}\label{eq:product-isometry}
 (\mathcal F,H)=(G,h_G)\otimes(V,H_V),\qquad
 V=H^q(Z,\Omega_Z^p\otimes M).
\end{equation}
If $V=0$ there is nothing to prove.  Otherwise choose
$0\ne\alpha\in\mathcal H^{p,q}(Z,M)$.

The product is complete and admits a positive line bundle, obtained
by externally tensoring a positive line on $Y$ with an ample line on
$Z$.  Proposition~\ref{prop:smooth-test} therefore gives
\[
 C_L=[i\Theta_{L,h},\Lambda_{\omega_X}]_{a+p,b+q}\ge0.
\]
The product cotangent splitting is orthogonal, and
$dV_X=dV_Y\,dV_Z$.  The curvature and Lefschetz contractions satisfy
\[
 i\Theta_{L,h}=\operatorname{pr}_Y^*i\Theta_{G,h_G}
                +\operatorname{pr}_Z^*i\Theta_{M,h_M},\qquad
 \Lambda_X=\Lambda_Y+\Lambda_Z,
\]
\[
 [i\Theta_G\wedge,\Lambda_Z]=0,\qquad
 [i\Theta_M\wedge,\Lambda_Y]=0.
\]
Consequently the component
$\Lambda^{a,b}T_Y^*\otimes\Lambda^{p,q}T_Z^*\otimes L$
is invariant, and its curvature operator is
\[
 C_L=C_G\otimes\Id+\Id\otimes C_M,\quad
 C_G=[i\Theta_G,\Lambda_Y]_{a,b},\quad
 C_M=[i\Theta_M,\Lambda_Z]_{p,q}.
\]
Fix $y$ and $\xi\in\Lambda^{a,b}T_y^*Y\otimes G_y$.  Pointwise
nonnegativity of $C_L$, followed by integration over the fiber,
gives
\[
 0\le\int_Z\langle C_L(\xi\otimes\alpha),
                              \xi\otimes\alpha\rangle\,dV_Z
   =\langle C_G\xi,\xi\rangle\|\alpha\|^2
      +|\xi|^2\int_Z\langle C_M\alpha,\alpha\rangle\,dV_Z .
\]
Since $\bar\partial_M\alpha=\bar\partial_M^*\alpha=0$, the
compact-fiber Bochner identity reads
\[
 0=\|\bar\partial_M\alpha\|^2+\|\bar\partial_M^*\alpha\|^2
   =\|D'_M\alpha\|^2+\|(D'_M)^*\alpha\|^2
         +\int_Z\langle C_M\alpha,\alpha\rangle\,dV_Z .
\]
In particular,
\[
 \int_Z\langle C_M\alpha,\alpha\rangle\,dV_Z
 =-\|D'_M\alpha\|^2-\|(D'_M)^*\alpha\|^2.
\]
Consequently,
\begin{equation}\label{eq:product-curvature-proof}
 \langle C_G\xi,\xi\rangle\|\alpha\|^2
 \ge |\xi|^2\bigl(\|D'_M\alpha\|^2+\|(D'_M)^*\alpha\|^2\bigr)
 \ge0.
\end{equation}
Dividing \eqref{eq:product-curvature-proof} by $\|\alpha\|^2>0$
shows that $C_G\ge0$ at every $y$.  The constant matrix $H_V$ in
the Gram formula implies
$i\Theta_{\mathcal F,H}=i\Theta_{G,h_G}\otimes\Id_V$.
Hence
$C_{\mathcal F;a,b}=C_G\otimes\Id_V\ge0$.
For any positive auxiliary line $A$ on $Y$, the total curvature
operator of $\mathcal F\otimes A$ is
\[
 C_{\mathcal F\otimes A;a,b}=C_{\mathcal F;a,b}+B_A\ge B_A>0.
\]
If $f$ is a datum in Definition~\ref{def:smooth-optimal}, its compact
support and smoothness imply $f\in L^2$ and finite inverse energy.
Lemma~\ref{lem:complete-estimate} and \eqref{eq:inverse-order} give
\[
 \bar\partial u=f,\qquad
 \|u\|^2\le\int_Yq_{C_{\mathcal F;a,b}+B_A}(f)\,dV_Y
       \le\int_Y\langle B_A^{-1}f,f\rangle\,dV_Y .
\]
\end{proof}

\begin{remark}\label{rem:product-attribution}
If $C_L\ge0$ is assumed directly, smooth compact
K\"ahler fibers suffice; projectivity was used only to provide a
positive auxiliary line on $X$.
\end{remark}

\begin{proposition}[Nonsplit metrics in range (ii)]\label{thm:variable-metric}
Keep the product geometry and holomorphic line bundle
$L=G\boxtimes M$ of Theorem~\ref{thm:product-descent}, but give $L$
an arbitrary smooth metric $h$.  In range (ii), its optimal
$(a+p,m+n)$ estimate implies the optimal $(a,m)$ estimate for the
prescribed harmonic metric on $R^n\pi_*(\Omega_{X/Y}^p\otimes L)$.
\end{proposition}
\begin{proof}
Let $V=H^n(Z,\Omega_Z^p\otimes M)$.  If $V=0$ the assertion is
empty.  Product base change identifies the direct image
holomorphically with $G\otimes V$.  Choose a basis $v_\nu$ and fixed smooth
representatives $\eta_\nu$.  On a base coordinate neighborhood
$U$, choose a holomorphic frame of $G$; the metric $h$ then gives
a smooth family of metrics $h_y$ on $M$.
Let $\Delta_y$ be the fiberwise $\bar\partial$-Laplacian in degree
$(p,n)$.  Its kernel has the constant dimension $\dim V$.
On fixed Sobolev spaces, the coefficients of $\Delta_y$ depend
smoothly on $y$.  Choose $\delta>0$ with
$\operatorname{Spec}(\Delta_{y_0})\cap\{|\zeta|\le2\delta\}=\{0\}$
at a fixed $y_0$.  After shrinking $U$, the circle
$\Gamma=\{|\zeta|=\delta\}$ remains in the resolvent set and
\[
 P_y=\frac1{2\pi i}\int_\Gamma(\zeta-\Delta_y)^{-1}\,d\zeta
\]
is the harmonic projection.  Indeed, the spectral projection has
locally constant rank, equal at $y_0$ to $\dim V$, while
$\dim\ker\Delta_y=\dim V$ for every $y$.
The resolvent depends smoothly on $y$ between Sobolev spaces:
the inverse identity
\[
 (\zeta-\Delta_y)^{-1}
 =(\zeta-\Delta_{y_0})^{-1}
 \bigl[\Id-(\Delta_y-\Delta_{y_0})
                   (\zeta-\Delta_{y_0})^{-1}\bigr]^{-1}
\]
is a locally convergent operator Neumann series, uniformly on
$\Gamma$.  Elliptic regularity and Sobolev embedding give
\[
 y\longmapsto P_y\eta_\nu\in C^\infty(U,H^s(Z))
 \quad\text{for every }s,\qquad
 (y,z)\longmapsto(P_y\eta_\nu)(z)\in C^\infty(U\times Z).
\]

Set $\alpha_\nu(y)=P_y\eta_\nu$.  If $G_y$ is the Green operator,
the top-antiholomorphic-degree Hodge identity gives
\[
 \eta_\nu-\alpha_\nu(y)
   =\Delta_yG_y\eta_\nu
   =\bar\partial_Z\bar\partial_{h_y}^*G_y\eta_\nu .
\]
Thus $\alpha_\nu(y)$ represents $v_\nu$, and
\[
 H_{\nu\bar\mu}(y)
   =\int_Z\langle\alpha_\nu(y),\alpha_\mu(y)\rangle_{h_y,\omega_Z}
                  \,dV_Z
\]
is the smooth Gram matrix of the prescribed metric, in the frame
of $G$ just chosen.  A change of that frame rescales $h_y$ by a
constant on each fiber and preserves its harmonic subspace.

For $0\le j\le m$, define the harmonic lift and the cohomology
projection by
\[
 J_jv=\sum_\nu\pi^*v_\nu\wedge\alpha_\nu(y),\qquad
 \Pi_jW=[W_{(a,j),(p,n)}]_{\bar\partial_Z}.
\]
Here $W$ has total bidegree $(a+p,j+n)$, and the subscript selects
horizontal degree $(a,j)$ and vertical degree $(p,n)$.
Every vertical $(p,n)$-form is $\bar\partial_Z$-closed.  In a
horizontal orthonormal form basis, apply the minimum-norm property
separately to each vertical coefficient $w_{I,J}$.  This gives
\[
 |\Pi_jW(y)|_H^2
   =\sum_{I,J}\|[w_{I,J}]\|_{H_y}^2
   \le\sum_{I,J}\|w_{I,J}\|_{L^2(Z,h_y)}^2
   \le\int_Z|W(y,z)|_{h,\omega_X}^2\,dV_Z .
\]
For a harmonic lift the first inequality is an equality and the
other product components vanish.  Integrating over $Y$ yields
\begin{equation}\label{eq:top-projection}
 \Pi_jJ_j=\Id,\qquad \|J_jv\|_X=\|v\|_Y,\qquad
 \|\Pi_jW\|_Y\le\|W\|_X.
\end{equation}
These identities remain valid after twisting by a base line $A$,
since its norm is constant on each fiber.

To compute the projection, choose the fixed Serre-dual basis
$\beta^\nu\in H^0(Z,\Omega_Z^{n-p}\otimes M^{-1})$, normalized by
\[
 \int_Z\eta_\mu\wedge\beta^\nu=\delta_\mu^\nu.
\]
In local holomorphic frames, write the selected component as
$\sum_{I,J}dt^I\wedge d\bar t^J\wedge w_{I,J}$.  Then
\[
 \Pi_jW
 =\sum_{I,J,\nu}
       \left(\int_Zw_{I,J}(y,z)\wedge\beta^\nu(z)\right)
                      dt^I\wedge d\bar t^J\otimes v_\nu.
\]
Thus $\Pi_j$ is independent of the coefficient metric.  For
$0\le j<m$, the component formula is
\[
 (\bar\partial_XW)_{(a,j+1),(p,n)}
 =\bar\partial_YW_{(a,j),(p,n)}
   +(-1)^{a+j+1}\bar\partial_ZW_{(a,j+1),(p,n-1)}
\]
For a smooth vertical $(p,n-1)$ coefficient $\gamma$,
\[
 \int_Z\bar\partial_Z\gamma\wedge\beta^\nu
  =\int_Z\bar\partial_Z(\gamma\wedge\beta^\nu)=0,
 \qquad \bar\partial_Z\beta^\nu=0.
\]
Consequently the coefficients of the preceding component formula
satisfy
\[
 \int_Z(\bar\partial_XW)_{(a,j+1),(p,n)}\wedge\beta^\nu
  =\bar\partial_Y
        \int_ZW_{(a,j),(p,n)}\wedge\beta^\nu.
\]
Equivalently,
\begin{equation}\label{eq:top-chain}
 \Pi_{j+1}\bar\partial_XW=\bar\partial_Y\Pi_jW.
\end{equation}
The same identity holds for $W$ in the maximal domain.  For every
compactly supported base test form $\tau$, the form
$\pi^*\tau\wedge\beta^\nu$ has compact support on $U\times Z$.
Distributional differentiation and fiber integration give
\[
 \operatorname{pr}_{Y*}
       \bigl(\bar\partial_Z\gamma\wedge\beta^\nu\bigr)=0,\qquad
 \operatorname{pr}_{Y*}
       \bigl(\bar\partial_Y\eta\wedge\beta^\nu\bigr)
   =\bar\partial_Y\operatorname{pr}_{Y*}
       \bigl(\eta\wedge\beta^\nu\bigr)
\]
when paired with $\tau$, for any distributional vertical
coefficients $\gamma,\eta$ of the indicated degrees.  These are
the same integration-by-parts identities, since
$\bar\partial_Z(\pi^*\tau\wedge\beta^\nu)=0$.
In particular neither summand in the component formula has to be
separately $L^2$.  The bounded projection estimate ensures that
both sides of \eqref{eq:top-chain}, whenever the left side is
defined from $\bar\partial_XW\in L^2$, represent the claimed
distributional derivative of an $L^2$ form.

Let $f$ be a compactly supported smooth $(a,m)$ datum and let $A>0$.
The lift $F=J_mf$ is smooth and supported in
$\operatorname{supp}f\times Z$, hence is $L^2$ with compact support.
It has total antiholomorphic degree $m+n$, so
$\bar\partial_XF=0$ for degree reasons.
For
\[
 B_X=[i\pi^*\Theta_A,\Lambda_X]_{a+p,m+n},\qquad
 B_Y=[i\Theta_A,\Lambda_Y]_{a,m},
\]
Choose horizontal unitary covectors with
$i\Theta_A=i\sum_{i=1}^m\lambda_i\,dt_i\wedge d\bar t_i$,
$\lambda_i>0$, and vertical unitary covectors on $Z$.
The pulled-back curvature has eigenvalues
$\lambda_1,\ldots,\lambda_m,0,\ldots,0$.  On a total
$(a+p,m+n)$ component with horizontal holomorphic index set $I$
and vertical holomorphic index set $K$,
formula~\eqref{eq:commutator-eigenvalue} gives
\[
 B_X\big|_{I,K}=\left(\sum_{i\in I}\lambda_i\right)\Id,
 \qquad |I|+|K|=a+p.
\]
Therefore $B_X\ge0$; its possible zero eigenspaces have $I=\varnothing$.
The lift $J_mf$ lies in the block $|I|=a\ge1$, $|K|=p$, on which
the eigenvalue is exactly that of $B_Y$ on the corresponding
horizontal component.  Hence
\[
 B_XJ_m=J_mB_Y,\qquad
 J_mf\perp\ker B_X,\qquad
 B_X^\dagger J_mf=J_m(B_Y^{-1}f).
\]
The lift isometry and its polarized identity now give the exact
inverse-energy formula
\[
 \int_Xq_{B_X}(J_mf)\,dV_X
 =\int_X\langle J_m(B_Y^{-1}f),J_mf\rangle\,dV_X
 =\int_Y\langle B_Y^{-1}f,f\rangle\,dV_Y.
\]
The hypothesis and Proposition~\ref{prop:smooth-test} give
$C_L\ge0$.  The curvature of $L\otimes\pi^*A$ is represented by
$C_L+B_X\ge B_X$.  Equations
\eqref{eq:extended-inverse}--\eqref{eq:inverse-order} give
\[
 \int_Xq_{C_L+B_X}(J_mf)\,dV_X
 \le\int_Xq_{B_X}(J_mf)\,dV_X<\infty.
\]
The complete-manifold estimate supplies $U$ with
$\bar\partial_XU=J_mf$ and
$\|U\|^2\le\int_Y\langle B_Y^{-1}f,f\rangle\,dV_Y$.
Set $u=\Pi_{m-1}U$.  Finally,
\[
 \bar\partial_Yu
   =\Pi_m\bar\partial_XU=\Pi_mJ_mf=f,\qquad
 \|u\|^2\le\|U\|^2
       \le\int_Y\langle B_Y^{-1}f,f\rangle\,dV_Y,
\]
by \eqref{eq:top-projection}--\eqref{eq:top-chain}.
\end{proof}

\begin{corollary}[Descent from a dense metrically split locus]
\label{cor:dense-product-descent}
Let $\pi:X\to Y$ be a proper holomorphic submersion, with $X$ and
$Y$ complete K\"ahler and admitting positive line bundles.  Assume
that the dimensions $h^q(X_y,\Omega_{X_y}^p\otimes L_y)$ are constant.
Suppose that, on a dense open subset $Y^\circ$, there are local
product charts in which
\[
 \omega_X=\pi^*\omega_Y+\omega_Z,\qquad
 (L,h)=(G_U,h_{G_U})\boxtimes(M_U,h_{M_U}),
\]
with fixed fiber data.  In either range \eqref{eq:ranges}, the
upstairs optimal estimate implies the downstairs optimal estimate.
\end{corollary}
\begin{proof}
The upstairs hypothesis and Proposition~\ref{prop:smooth-test}
give $C_L\ge0$ on $X$.  Fix a product chart over $U\subset Y^\circ$.
If the common cohomology rank is zero, the conclusion is immediate.
Otherwise choose a nonzero harmonic fiber class $\alpha$.  For a horizontal tensor
$\xi$, the same pointwise splitting and compact-fiber integration
used in \eqref{eq:product-curvature-proof} give
\[
 \langle C_{G_U;a,b}\xi,\xi\rangle\|\alpha\|^2
 \ge|\xi|^2\bigl(\|D'_{M_U}\alpha\|^2+
                         \|(D'_{M_U})^*\alpha\|^2\bigr)\ge0.
\]
The local Gram matrix is $h_{G_U}H_{V_U}$ with $H_{V_U}$ constant,
so
\[
 C_{\mathcal F;a,b}
   =[i\Theta_{\mathcal F,H},\Lambda_Y]_{a,b}
   =C_{G_U;a,b}\otimes\Id_{V_U}\ge0\quad\text{on }U.
\]
To pass across $Y\setminus Y^\circ$, constant cohomology dimension
and Grauert's theorem give local freeness and base change for
$\mathcal F$.  In a local smooth trivialization of the proper
submersion, the fiber Laplacians are smooth elliptic families.
Their kernels have constant dimension; the spectral projection
construction above therefore gives smooth harmonic representatives
of local frames.  Their Gram entries are integrals over a fixed
compact smooth fiber of smoothly varying coefficients, so
$H_{\nu\bar\mu}\in C^\infty(Y)$.  In particular
$C_{\mathcal F;a,b}$ is continuous on all of $Y$.

Fix $y_0\in Y$, a sequence $y_\ell\in Y^\circ$ with $y_\ell\to y_0$,
and $\xi_0\in\Lambda^{a,b}T_{y_0}^*Y\otimes\mathcal F_{y_0}$.
Choose a local smooth section $\xi$ with $\xi(y_0)=\xi_0$.  Then
\[
 \langle C_{\mathcal F;a,b}(y_0)\xi_0,\xi_0\rangle
 =\lim_{\ell\to\infty}
   \langle C_{\mathcal F;a,b}(y_\ell)\xi(y_\ell),\xi(y_\ell)\rangle
 \ge0 .
\]
Thus $C_{\mathcal F;a,b}\ge0$ everywhere.  For each positive
auxiliary line $A$ on the complete base,
$C_{\mathcal F\otimes A;a,b}\ge B_A>0$.
Lemma~\ref{lem:complete-estimate} and
\eqref{eq:inverse-order} give the claimed optimal estimate.
\end{proof}
Infinitesimal triviality supplies dense holomorphic product charts,
but not the two metric splittings in this corollary.  The
counterexample in Corollary~\ref{cor:counter-higher-q} is already a
globally trivial family.


\section{Top direct images and the smooth adjoint case}\label{sec:smooth-duality}

Let $\pi:X\to Y$ be a proper surjective holomorphic submersion of
relative dimension $n$ between K\"ahler manifolds.  A smooth Hermitian
bundle $(E,h)$ is \emph{dual Nakano semipositive} if $(E^*,h^*)$ is
Nakano seminegative.  The submersion is \emph{infinitesimally trivial}
if its Kodaira--Spencer maps
\[
 \rho_y:T_y^{1,0}Y\longrightarrow H^1(X_y,T_{X_y})
\]
vanish for every $y$.  In this section the dimension
$h^n(X_y,E_y)$ is assumed locally constant.


Theorem~\ref{thm:top-dual} is the metric Serre-dual form of the negative-adjoint
mechanism in Liu--Yang \cite[Theorem 4.1]{LY14}.  We give the metric
identification and the dense-locus argument explicitly.

\begin{lemma}[Isometric Serre duality]\label{lem:serre-isometry}
For a compact complex $n$-fold $Z$ and a smooth Hermitian bundle
$(E,h)$, the quotient norm on $H^n(Z,E)$ is independent of the
Hermitian metric on $Z$.  With the canonical $L^2$ norm on
$H^0(Z,K_Z\otimes E^*)$, Serre duality is an isometry.  In the family,
\begin{equation}\label{eq:serre-isometry}
 (R^n\pi_*E,H)
 \simeq\bigl(\pi_*(K_{X/Y}\otimes E^*),G\bigr)^*
\end{equation}
as smooth Hermitian holomorphic bundles.
\end{lemma}
\begin{proof}
\emph{The ambient Hilbert norm.}
In a holomorphic coordinate chart put
$d\lambda_z=i^{n^2}dz_1\wedge\cdots\wedge dz_n
\wedge d\bar z_1\wedge\cdots\wedge d\bar z_n$.
For $a=a_0\,d\bar z_1\wedge\cdots\wedge d\bar z_n$ and
$\omega=i\sum g_{j\bar k}dz_j\wedge d\bar z_k$,
\[
 |a|_{h,\omega}^2=(\det g)^{-1}|a_0|_h^2,
 \qquad dV_\omega=(\det g)d\lambda_z,
 \qquad |a|_{h,\omega}^2dV_\omega=|a_0|_h^2d\lambda_z.
\]
The same identity, with $h^*$ in place of $h$, holds for
$s=s_0\,dz_1\wedge\cdots\wedge dz_n$.
Thus both Hilbert norms are independent of $\omega$.
The subspace
\[
 \mathcal B=\overline{\bar\partial A^{0,n-1}(Z,E)}^{\,L^2_{0,n}}
\]
is the closure of the same set in this fixed Hilbert norm.
For a smooth representative $a$, compact Hodge theory gives
$a=a_h+\bar\partial b$ with $a_h\perp\mathcal B$, and therefore
\[
 \inf_{c\in\mathcal B}\|a+c\|^2=\|a_h\|^2,
 \qquad \|a_h+c\|^2=\|a_h\|^2+\|c\|^2.
\]
This proves the asserted independence of the quotient norm.

\emph{The duality map and equality of norms.}
The exterior-product pairing defines an antilinear isometry $J$ by
\[
 s\wedge a=\langle s,Ja\rangle_{h^*,\omega}dV_\omega,
 \qquad
 J:\Lambda^{0,n}T_Z^*\otimes E\longrightarrow
       \Lambda^{n,0}T_Z^*\otimes E^*.
\]
The fixed unit phase from the exterior ordering is included in $J$.
With the Hermitian products linear in their first entry, one has
\[
 (Ja,Jc)_{L^2}=(c,a)_{L^2},\qquad
 \int_ZJa\wedge\bar\partial b=(\bar\partial b,a)_{L^2}.
\]
For smooth $a$ and $b\in A^{0,n-1}(Z,E)$, Stokes gives
\[
 \int_ZJa\wedge\bar\partial b
 =(-1)^{n+1}\int_Z\bar\partial(Ja)\wedge b.
\]
Since $b$ is arbitrary, the nondegenerate test-form pairing implies
\[
 a\perp\bar\partial A^{0,n-1}(Z,E)
 \ \Longleftrightarrow\
 \int_Z Ja\wedge\bar\partial b=0\quad\text{for all }b
 \ \Longleftrightarrow\ \bar\partial(Ja)=0.
\]
Every $(0,n)$-form is $\bar\partial$-closed.  Consequently $J$
identifies harmonic $(0,n)$-forms antiunitarily with
$H^0(Z,K_Z\otimes E^*)$.  For the minimum representative $a_h$,
the integral $\int_Zs\wedge a$ depends only on $[a]$, because
$\bar\partial s=0$.  Moreover,
\[
 \left|\int_Zs\wedge a\right|
 =\left|\int_Zs\wedge a_h\right|
 \le\|s\|_{L^2}\|a_h\|_{L^2},\qquad
 \int_ZJa_h\wedge a_h=\|a_h\|_{L^2}^2.
\]
The upper bound is thus attained, and
\begin{equation}\label{eq:serre-norm-sup}
 \|[a]\|_H^2
 =\sup_{s\ne0}
   \frac{|\int_Zs\wedge a|^2}{\|s\|_G^2};
\end{equation}
For the zero class both sides are zero; otherwise use $s=Ja_h$.

\emph{Holomorphic and smooth dependence on the base.}
Constant fiber dimension and Grauert's theorem \cite{Gra60} give local freeness
and base change.  The Serre pairing is holomorphic by relative duality
\cite{RRV71}: for a relative
Dolbeault representative, its base $\bar\partial$ derivative differs
from the representative of the cohomological derivative by a
fiberwise exact term, whose integral against a holomorphic top form
vanishes.  Let $s_1,\ldots,s_r$ be a local holomorphic frame of
$\pi_*(K_{X/Y}\otimes E^*)$ and let $e_1,\ldots,e_r$ be its
Serre-dual frame.  Fiberwise Serre duality gives the constant rank
of the first bundle, and
\[
 G_{\alpha\bar\beta}(y)
 =\int_{X_y}\langle s_\alpha(y),s_\beta(y)\rangle_{h^*,\omega_y}
                 dV_{\omega_y},\qquad
 H(y)=(G(y)^{-1})^{\mathsf T}.
\]
To justify smoothness of $G$, use a local $C^\infty$ trivialization
$\pi^{-1}(U)\simeq U\times X_{y_0}$ and a finite partition of unity
on its compact fibers.  Each matrix entry is then a finite sum
$\int_{X_{y_0}}a_{\alpha\bar\beta}(y,z)\,dV_0(z)$ with smooth
integrand.  On relatively compact $U'\Subset U$, every derivative
of the integrand is bounded on $\overline{U'}\times X_{y_0}$, so
\[
 \partial_y^\gamma G_{\alpha\bar\beta}(y)
 =\int_{X_{y_0}}\partial_y^\gamma
       a_{\alpha\bar\beta}(y,z)\,dV_0(z).
\]
The matrix $G(y)$ is positive definite, hence its inverse and $H$
are smooth.  Formula~\eqref{eq:serre-norm-sup} identifies this dual
metric with the prescribed harmonic metric.
\end{proof}

\begin{lemma}[Dense product locus]\label{lem:dense-product}
An infinitesimally trivial proper holomorphic submersion is
holomorphically locally a product over a dense open $Y^\circ\subset Y$.
Locally on $Y$, the complement can be chosen analytic.
\end{lemma}
\begin{proof}
Work over a connected coordinate neighborhood and remove the analytic
locus where $h^1(X_y,T_{X_y})$ exceeds its generic value.  Grauert's
theorem gives, on the remaining dense open set,
\[
 \mathcal T=R^1\pi_*T_{X/Y}\ \text{locally free},\qquad
 \mathcal T\otimes\mathbb C(y)\simeq H^1(X_y,T_{X_y}).
\]
These local dense open sets cover a dense open $Y^\circ\subset Y$.
Apply $\pi_*$ to
\[
 0\longrightarrow T_{X/Y}\longrightarrow T_X
   \xrightarrow{d\pi}\pi^*T_Y\longrightarrow0.
\]
The resulting exact sequence contains
\[
 \pi_*T_X\longrightarrow\pi_*\pi^*T_Y
       \xrightarrow{\delta}\mathcal T.
\]
For $\xi_j=\partial/\partial t_j$ regarded as a section of
$\pi_*\pi^*T_Y$, base change identifies
\[
 \delta(\xi_j)(y)=\rho_y(\xi_j(y))=0.
\]
Write $\delta(\xi_j)=\sum_\alpha f_{j\alpha}e_\alpha$ in a local
frame of $\mathcal T$.  Thus $f_{j\alpha}(y)=0$ for every $y$, and
$\delta(\xi_j)=0$ as a section.  Exactness, followed by one shrinking
for the finitely many indices $j$, gives
\[
 V_j\in H^0(\pi^{-1}(U),T_X),\qquad d\pi(V_j)=\xi_j.
\]

Center the coordinates at $y_0=0$ and let $\Phi_j^\zeta$ be the
local holomorphic flow of $V_j$.  Compactness of $X_0$ gives a
common $\varepsilon>0$ on which all successive flows below are
defined near this fiber for $|t_j|<\varepsilon$.  Their base
coordinates satisfy
\[
 \frac{d}{d\zeta}\pi(\Phi_j^\zeta(x))=e_j,
 \qquad \pi(\Phi_j^\zeta(x))=\pi(x)+\zeta e_j.
\]
Consequently
\[
 F(t,z)=\Phi_m^{t_m}\circ\cdots\circ\Phi_1^{t_1}(z),
 \qquad F:\Delta_\varepsilon^m\times X_0\longrightarrow X,
 \qquad \pi F(t,z)=t.
\]
At $(0,z)$, $dF$ identifies the vertical tangent with $T_zX_0$
and sends the coordinate directions to $V_j(z)$, so it is invertible.
Compactness allows a common shrinking on which it stays invertible.
For fixed $t$, the composition of flows is injective; different $t$
have different images under $\pi$.  Thus $F$ is a biholomorphism
onto an open $W$ containing $X_0$, with inverse
\[
 x\longmapsto\left(t(x),
   \Phi_1^{-t_1(x)}\circ\cdots\circ\Phi_m^{-t_m(x)}(x)\right).
\]
These two maps are inverse.  The relatively closed set
$C=\pi^{-1}(\Delta_\varepsilon^m)\setminus W$ has closed image
$\pi(C)$ by properness, and $0\notin\pi(C)$.
Choose $U'\ni0$ disjoint from $\pi(C)$.  Then
\[
 \pi^{-1}(U')\subset W,\qquad
 F|_{U'\times X_0}:U'\times X_0\xrightarrow{\sim}\pi^{-1}(U').
\]
This proves local product structure at every point of $Y^\circ$.
\end{proof}

\begin{proof}[Proof of Theorem~\ref{thm:top-dual}]
Put $V=E^*$.  By Lemma~\ref{lem:dense-product}, it suffices first
to work on a product chart $U\times Z$ over $Y^\circ$.
Let $P_y$ be the
orthogonal projection onto $H^0(Z,K_Z\otimes V_y)$ and let
$\nabla_i,\nabla_{\bar j}$ be the horizontal Chern derivatives of $V$.
The canonical density of Lemma~\ref{lem:serre-isometry} permits
integration with a fixed fiber metric.  Therefore the metric
derivative contains no derivative of a fiber volume form:
\[
 \partial_i(s,t)_{L^2}
   =(\nabla_i s,t)_{L^2}+(s,\nabla_{\bar i}t)_{L^2}.
\]
For $s,t$ in the direct image, orthogonality gives
$(\nabla_i s,t)=(P\nabla_i s,t)$.
Moreover, a holomorphic direct-image section is a holomorphic
relative top form on $U\times Z$, so $\nabla_{\bar j}s=0$.
Thus the projected connection has the correct metric compatibility
and $(0,1)$ part, and is the Chern connection:
\[
 D_i^G=P\nabla_i,\qquad B_i=(1-P)\nabla_i|_{\operatorname{im}P}.
\]
The operator $B_i$ is tensorial: for a holomorphic function $f$ on $U$,
\[
 (1-P)\nabla_i(fs)
 =(1-P)((\partial_i f)s+f\nabla_i s)=fB_i s.
\]
We now compute its curvature contribution rather than assume a
subbundle formula.  The result is
\begin{equation}\label{eq:negative-adjoint-curvature}
 \langle R^G_{i\bar j}s,t\rangle_G
 =\int_Z\langle R^V_{i\bar j}s,t\rangle_{h^*,\omega_Z}\,dV_Z
       -(B_i s,B_j t)_{L^2}.
\end{equation}
Choose a local holomorphic frame normal for $G$ at $y_0$, and take
$s,t$ with constant coefficients in this frame.  The normality
conditions are $G(y_0)=\Id$ and $\partial_iG(y_0)=0$, hence
$P\nabla_i s(y_0)=0$.  Since
$R^V_{i\bar j}=[\nabla_i,\nabla_{\bar j}]$ and
$\nabla_{\bar j}s=0$, differentiating the preceding metric identity gives
\[
 \begin{aligned}
 \bar\partial_j\partial_i(s,t)_{L^2}
   &=(\nabla_{\bar j}\nabla_i s,t)_{L^2}
          +(\nabla_i s,\nabla_j t)_{L^2}\\
   &=-\int_Z\langle R^V_{i\bar j}s,t\rangle_{h^*,\omega_Z}\,dV_Z
          +(\nabla_i s,\nabla_j t)_{L^2}.
 \end{aligned}
\]
At $y_0$, $\nabla_i s=B_i s$ and
$\langle R^G_{i\bar j}s,t\rangle
=-\bar\partial_j\partial_i(s,t)_{L^2}$.
Substitution proves
\eqref{eq:negative-adjoint-curvature}.
For arbitrary $s_1,\ldots,s_m$,
\[
 \sum_{i,j}\langle R^G_{i\bar j}s_i,s_j\rangle_G
 =\int_Z\sum_{i,j}\langle R^V_{i\bar j}s_i,s_j\rangle_{h^*,\omega_Z}\,dV_Z
   -\Bigl\|\sum_iB_i s_i\Bigr\|^2\le0.
\]
For the first term, write $s_i(z)=v_i(z)\,dz_1\wedge\cdots\wedge dz_n$
in a fiber coordinate chart.  Nakano seminegativity of $V$ applied
to $\sum_i\partial_{t_i}\otimes v_i(z)$ gives
\[
 \sum_{i,j}\langle R^V_{i\bar j}v_i(z),v_j(z)\rangle_{h^*}\le0.
\]
Multiplication by the positive canonical density and integration
prove the asserted sign.  Thus $G$ is Nakano seminegative on
$Y^\circ$.

To pass across its complement, choose any $y_0\in Y$ and extend
the entries of a tensor in $T_{y_0}^{1,0}Y\otimes G_{y_0}$ smoothly
to $s_1(y),\ldots,s_m(y)$ in local frames.  Lemma~\ref{lem:serre-isometry}
makes the function
\[
 Q(y)=\sum_{i,j}\langle R^G_{i\bar j}(y)s_i(y),s_j(y)\rangle_G
\]
continuous.  For $y_\nu\in Y^\circ$ tending to $y_0$,
$Q(y_0)=\lim_\nu Q(y_\nu)\le0$.
The isometry \eqref{eq:serre-isometry} gives the asserted dual
Nakano semipositivity of $H$.
\end{proof}

\begin{remark}[Smooth semipositive adjoints and the prescribed metric]
\label{rem:mt-comparison}
Let $\pi:X\to Y$ be a proper K\"ahler holomorphic submersion and
let $(L,h)$ be a smooth semipositive line bundle on $X$.  For every
$0\le q\le n$, the ordinary fiberwise harmonic metric on
\[
 \mathcal E_q=R^q\pi_*(K_{X/Y}\otimes L)
\]
is Nakano semipositive.  This is a consequence of
Mourougane--Takayama \cite[Theorem 1.1]{MT08}, with the following
explicit comparison of metrics.  Connected fibers may be assumed:
locally on the base, disconnected fibers give an orthogonal direct
sum of the connected cases.

Write $m=\dim Y$, $c_d=i^{d^2}$ and
$\omega_y=\omega_X|_{X_y}$.  In the notation of
\cite[Corollary 4.4]{MT08}, take the relative injection
\[
 S_\omega:\mathcal E_q\longrightarrow
        \pi_*(\Omega_{X/Y}^{n-q}\otimes L).
\]
For a local holomorphic section $s$, put $\sigma=S_\omega s$.
The base-change assertion in \cite[Lemma 4.9]{MT08} is
\[
 \mathcal E_q\otimes\mathbb C(y)
    \xrightarrow{\sim}H^q(X_y,K_{X_y}\otimes L_y).
\]
Thus every fiber class is the value of a local holomorphic section.
Corollary 4.4 and Lemma 4.7(1) of \cite{MT08} give
\begin{equation}\label{eq:mt-three-identities}
 \bar\partial_y\sigma_y=0,\qquad
 D'_{h_y}\sigma_y=0,\qquad
 s_y=\left[\frac{c_{n+m-q}}{q!}\omega_y^q\wedge\sigma_y\right].
\end{equation}
Work over a Stein coordinate ball on the base.  The Leray spectral
sequence and Cartan's theorem B identify sections of $\mathcal E_q$
with degree-$q$ cohomology on its inverse image.  If $\rho$ is a
Dolbeault representative of $s$, the identity
$(c_{n+m-q}/q!)L_\omega S_\omega=\Id$ says, after the relative
canonical bundle is identified using $dt=dt_1\wedge\cdots\wedge dt_m$,
\[
 (\rho-c_{n+m-q}\omega_X^q\wedge\sigma/q!)\wedge dt
       =\bar\partial(b\wedge dt).
\]
Restricting the relative coefficients to $X_y$ kills all horizontal
differentials and gives
\[
 \rho_y-c_{n+m-q}\omega_y^q\wedge\sigma_y/q!
      =\bar\partial_y b_y.
\]

Set $\alpha_y=c_{n+m-q}\omega_y^q\wedge\sigma_y/q!$.
The fiber star identity \cite[Corollary 2.2]{MT08} yields
\[
 *_y\alpha_y=\frac{c_{n+m-q}}{c_{n-q}}\sigma_y.
\]
Since $\omega_y$ is K\"ahler, \eqref{eq:mt-three-identities} gives
\[
 \bar\partial_y\alpha_y=0,\qquad
 \bar\partial_y^*\alpha_y
 =-*_yD'_{h_y}(*_y\alpha_y)=0.
\]
To check the minimizing property explicitly, every representative
of the same class has the form $\alpha_y+\bar\partial_y v$.  On the
compact fiber, integration by parts gives
\[
 (\alpha_y,\bar\partial_y v)_{L^2}
   =(\bar\partial_y^*\alpha_y,v)_{L^2}=0,\qquad
 \|\alpha_y+\bar\partial_yv\|^2
   =\|\alpha_y\|^2+\|\bar\partial_yv\|^2.
\]
The identity extends to the maximal domain by graph approximation.
Thus $\alpha_y$ is the ordinary harmonic, minimum-norm representative.
For the metric of \cite[Definition 5.1]{MT08}, the pointwise
Lefschetz identity gives
\[
 g_{\mathrm{MT},y}(s_y,s_y)
   =\frac{c_{n-q}}{q!}\int_{X_y}
      \omega_y^q\wedge\sigma_y\wedge h_y\overline{\sigma_y}
   =\int_{X_y}|\sigma_y|_{h_y,\omega_y}^2\,dV_y.
\]
The star is an isometry and $|c_{n+m-q}/c_{n-q}|=1$, so
\begin{equation}\label{eq:mt-metric-equality}
 H_y(s_y,s_y)=\|\alpha_y\|^2_{L^2}
             =\|\sigma_y\|^2_{L^2}
             =g_{\mathrm{MT},y}(s_y,s_y).
\end{equation}
This equality holds on every fiber vector.  Polarization therefore
gives $H_y(s_y,t_y)=g_{\mathrm{MT},y}(s_y,t_y)$ for all $s_y,t_y$.

Mourougane--Takayama's curvature theorem therefore applies to
the prescribed $H$.  For a complete K\"ahler base and
$1\le b\le m$, Nakano semipositivity implies
\[
 [i\Theta_{\mathcal E_q,H},\Lambda_{\omega_Y}]_{m,b}\ge0.
\]
Indeed, write in unitary coordinates
$u=\sum_{|J|=b}u_J\,dz_1\wedge\cdots\wedge dz_m\wedge d\bar z_J$,
and let $u_{iK}$ be its signed alternating coefficients, zero if
$i\in K$.  Then
\[
 \langle[i\Theta_{\mathcal E_q,H},\Lambda]u,u\rangle
   =\sum_{|K|=b-1}\sum_{i,j}
        \langle R_{i\bar j}u_{iK},u_{jK}\rangle_H\ge0.
\]
Each inner sum is the Nakano form evaluated on
$\sum_i\partial_{z_i}\otimes u_{iK}$.
After twisting by a positive auxiliary line $A$, its commutator is
\[
 C_{\mathcal E_q\otimes A;m,b}
   =[i\Theta_{\mathcal E_q,H},\Lambda]_{m,b}+B_A\ge B_A>0.
\]
For every datum in Definition~\ref{def:smooth-optimal},
Lemma~\ref{lem:complete-estimate} and inverse order
\eqref{eq:inverse-order} give
\[
 \bar\partial u=f,\qquad
 \|u\|^2_{Hh_A}
 \le\int_Y\langle B_A^{-1}f,f\rangle_{Hh_A}\,dV_Y,
 \quad B_A=[i\Theta_A\otimes\operatorname{Id},\Lambda_{\omega_Y}]_{m,b}.
\]
This proves general descent in range (i) for smooth semipositive
$L$, including $q>0$.  The same comparison works for a smooth Nakano
semipositive
vector coefficient.
\end{remark}


\section{Singular product metrics}\label{sec:singular-product}

\begin{lemma}[Scalar estimates in the two extremal ranges]
\label{lem:singular-scalar}
Let $(G,h_G)$ be a holomorphic line bundle with psh local weights.
If $(Y,\omega_Y)$ is complete K\"ahler of dimension $m$, then
$(G\otimes\mathcal I(h_G),h_G)$ satisfies the weighted optimal
$(m,b)$ estimate for $1\le b\le m$ and the weighted optimal
$(a,m)$ estimate for $1\le a\le m$.  The second assertion holds
on any K\"ahler manifold.
\end{lemma}
\begin{proof}
Fix a smooth positive auxiliary metric $h_A$ on a holomorphic line
bundle $A$, and write $E=G\otimes A$ and $h=h_Gh_A$.
For the first assertion, choose a smooth reference metric $h_{G,0}$
and define the globally defined function $\chi$ by
\[
 h_G=h_{G,0}e^{-\chi},\qquad h_0=h_{G,0}h_A,\qquad h=h_0e^{-\chi}.
\]
Locally $\chi$ is a psh function plus a smooth function.  Its
distributional complex Hessian is therefore a current of order zero,
whose singular part is nonnegative.  Denote by
\[
 R=i\Theta_{E,h_0}+(i\partial\bar\partial\chi)_{\rm ac}
   =i\Theta_{A,h_A}+(i\Theta_{G,h_G})_{\rm ac}
\]
the absolutely continuous curvature form.  It has measurable
coefficients and satisfies $R\ge i\Theta_{A,h_A}$ almost everywhere.
On $(m,b)$-forms, put
\[
 C=[R,\Lambda_Y]_{m,b},\qquad
 B_A=[i\Theta_{A,h_A},\Lambda_Y]_{m,b}.
\]
The eigenvalue formula \eqref{eq:commutator-eigenvalue} gives
$C\ge B_A>0$, and hence
\[
 \langle C^{-1}f,f\rangle_h
 \le\langle B_A^{-1}f,f\rangle_h
 \quad\text{almost everywhere}.
\]
If the right-hand side has finite integral, then on each compact
$K\subset Y$ there is $c_K>0$ with
\[
 B_A^{-1}\ge c_K\Id,\qquad
 \int_K|f|_h^2\,dV_Y\le c_K^{-1}I_A(f).
\]
In particular, $f$ has locally $L^2$ coefficients also for the smooth
reference metric, since $h$ is locally bounded below.

Apply \cite[Section~5, Theorem~5.1]{Dem82} to the smooth line bundle
$(E,h_0)$ and the singular weight $\chi$.  The input of that theorem
is precisely a distributionally closed $(m,b)$-form with locally
$L^2$ coefficients and finite inverse energy for the absolutely
continuous curvature.  Its conclusion, in the present notation, is
\[
 \bar\partial u=f,\qquad
 \int_Y|u|_h^2\,dV_Y
 \le\int_Y\langle C^{-1}f,f\rangle_h\,dV_Y
 \le I_A(f).
\]
This proves the first estimate for the specified $\omega_Y$.

For the second assertion, set $k=h^{-1}$ on $E^*$ and $r=m-a$.
The metric $k$ is positive and finite almost everywhere and locally
bounded above.  Define
\[
 B=[i\Theta_{A,h_A},\Lambda_Y]_{a,m},\qquad
 Q=[-i\Theta_{A,h_A},\Lambda_Y]_{r,0}.
\]
If $i\Theta_{A,h_A}$ is diagonal in a unitary coframe, with eigenvalues
$\lambda_1,\ldots,\lambda_m>0$, then
\[
 B(dz^I\wedge d\bar z^{1}\wedge\cdots\wedge d\bar z^{m})
   =\beta_I\,dz^I\wedge d\bar z^{1}\wedge\cdots\wedge d\bar z^{m},
 \qquad \beta_I=\sum_{i\in I}\lambda_i,\quad |I|=a,
\]
whereas $Q(dz^{I^c})=\beta_I dz^{I^c}$.  In particular, both
operators are positive.  We prove the compactly supported dual-test
inequality
\begin{equation}\label{eq:singular-dual-test}
 \int_Y\langle Qw,w\rangle_k\,dV_Y
 \le\|\bar\partial w\|_k^2,\qquad
 w\in C_c^\infty(Y,\Lambda^{r,0}T_Y^*\otimes E^*).
\end{equation}
All terms are finite because $k$ is locally bounded above.

Choose a coordinate ball with holomorphic frames and write
$h_G=e^{-\varphi}$ and $h_A=e^{-\psi}$, where $\varphi$ is psh and
$\psi$ is smooth.  On a smaller ball, convolution with radial
regularizing kernels gives smooth psh functions
$\varphi_\varepsilon\downarrow\varphi$.  For a compactly supported
smooth $(r,0)$-form $v$ in this ball, set
\[
 k_\varepsilon=e^{\varphi_\varepsilon+\psi},\qquad
 s_\varepsilon=e^{(\varphi_\varepsilon+\psi)/2}.
\]
The coefficient curvature is
\[
 i\Theta_{E^*,k_\varepsilon}
 =-i\partial\bar\partial(\varphi_\varepsilon+\psi).
\]
On $(r,0)$-forms, the commutator
$[-i\partial\bar\partial\varphi_\varepsilon,\Lambda_Y]$ is
nonnegative: its eigenvalues are the sums of the complementary
$m-r$ eigenvalues of $i\partial\bar\partial\varphi_\varepsilon$.
Also $\bar\partial^*v=0$ by bidegree.  The smooth
Bochner--Kodaira--Nakano identity therefore gives
\begin{equation}\label{eq:dual-local-smooth}
 \|\bar\partial v\|_{k_\varepsilon}^2
 \ge\|D'_{k_\varepsilon}v\|_{k_\varepsilon}^2
   +\|(D'_{k_\varepsilon})^*v\|_{k_\varepsilon}^2
   +\int\langle Qv,v\rangle_{k_\varepsilon}\,dV_Y.
\end{equation}
By the K\"ahler identity, $(D'_{k_\varepsilon})^*$ is
$i[\Lambda_Y,\bar\partial]$, up to the fixed conventional sign.
Write this operator as $P$.  It is independent of the coefficient
metric, so the second term on the right equals
$\|s_\varepsilon Pv\|^2$.

To pass to the limit in the connection term, put
$t_\varepsilon=e^{\varphi_\varepsilon/2}$.
These functions are nonnegative smooth psh functions and are uniformly
bounded on compact subballs.  For a real compactly supported cutoff
$\eta$, the Euclidean Laplacian satisfies
\[
 \Delta(t_\varepsilon^2)
 =2t_\varepsilon\Delta t_\varepsilon
       +2|\nabla t_\varepsilon|^2
 \ge2|\nabla t_\varepsilon|^2.
\]
Integration by parts consequently gives the uniform local bound
\[
 2\int\eta^2|\nabla t_\varepsilon|^2
 \le\int t_\varepsilon^2\Delta(\eta^2)
 \le\bigl(\sup_{\operatorname{supp}\eta}t_\varepsilon^2\bigr)
          \int|\Delta(\eta^2)|.
\]
Since $t_\varepsilon\to e^{\varphi/2}$ almost everywhere and the
functions are uniformly bounded, dominated convergence and weak
compactness give, after passage to a subsequence,
\[
 s_\varepsilon\longrightarrow s=e^{(\varphi+\psi)/2}
       \text{ strongly in }L^2_{\rm loc},\qquad
 s_\varepsilon\rightharpoonup s
       \text{ weakly in }W^{1,2}_{\rm loc}.
\]
More explicitly, on each compact subball $K$,
\[
 \int_K|s_\varepsilon-s|^2\longrightarrow0,\qquad
 \sup_\varepsilon\int_K|ds_\varepsilon|^2<\infty.
\]
For every compactly supported smooth test function $\zeta$ and every
real coordinate $x_j$, any weak derivative limit $a_j$ satisfies
\[
 \int a_j\zeta
 =\lim_\varepsilon\int (\partial_{x_j}s_\varepsilon)\zeta
 =-\lim_\varepsilon\int s_\varepsilon\partial_{x_j}\zeta
 =-\int s\partial_{x_j}\zeta.
\]
Thus $a_j=\partial_{x_j}s$ distributionally.
In the chosen holomorphic frame,
\[
 D'_{k_\varepsilon}v
   =\partial v+\partial\log k_\varepsilon\wedge v,\qquad
 s_\varepsilon D'_{k_\varepsilon}v
   =s_\varepsilon\partial v+2\partial s_\varepsilon\wedge v
   \rightharpoonup T_sv,
\]
where $T_sv=s\partial v+2\partial s\wedge v\in L^2_{\rm loc}$.
The remaining terms of \eqref{eq:dual-local-smooth} converge by
dominated convergence: $v$, $Pv$, $\bar\partial v$, and $Q$ are smooth
on a fixed compact support, and $k_\varepsilon$ is uniformly bounded
there.  In particular,
\[
 \|T_sv\|^2\le\liminf_{\varepsilon\to0}
       \|s_\varepsilon D'_{k_\varepsilon}v\|^2,\qquad
 \|s_\varepsilon Pv-sPv\|\longrightarrow0,\qquad
 \|s_\varepsilon\bar\partial v-s\bar\partial v\|\longrightarrow0.
\]
Passing to the limit in \eqref{eq:dual-local-smooth} therefore yields
\begin{equation}\label{eq:dual-local-limit}
 \|s\bar\partial v\|^2\ge\|T_sv\|^2+\|sPv\|^2
                          +\int\langle Qv,v\rangle_k\,dV_Y.
\end{equation}

The densities in this inequality are invariant under holomorphic
frame changes.  If the frame of $E^*$ is multiplied by a nowhere
vanishing holomorphic function $g$, then the coefficients and square
root of the metric become
\[
 v'=g^{-1}v,\qquad s'=|g|s.
\]
Since $2\partial|g|=(|g|/g)\partial g$, the weak product rule gives
\[
 T_{s'}v'=\frac{|g|}{g}T_sv,\qquad
 s'Pv'=\frac{|g|}{g}sPv,\qquad
 s'\bar\partial v'=\frac{|g|}{g}s\bar\partial v.
\]
Thus their squared norms agree on overlaps.

Cover $\operatorname{supp}w$ by finitely many such balls and choose
real smooth functions $\rho_j$ supported in them with
$\sum_j\rho_j^2=1$ on a neighborhood of $\operatorname{supp}w$.
Apply \eqref{eq:dual-local-limit} to $\rho_jw$ and sum.  The product
rules are
\[
 \bar\partial(\rho_jw)=\rho_j\bar\partial w
                             +\bar\partial\rho_j\wedge w,\qquad
 T_s(\rho_jw)=\rho_jT_sw+s\partial\rho_j\wedge w,
\]
and
\[
 P(\rho_jw)=\rho_jPw+[P,\rho_j]w,\qquad
 |[P,\rho_j]w|=|(\partial\rho_j\wedge)^*w|.
\]
The last identity follows directly from
$P=\pm i[\Lambda_Y,\bar\partial]$; its possible unit factor has no
effect on the norm.  Every cross term in the sum vanishes because
$\sum_j\rho_j\,d\rho_j=0$; explicitly,
\begin{align*}
 \sum_j\|s\bar\partial(\rho_jw)\|^2
 &=\|s\bar\partial w\|^2
        +\sum_j\|s\bar\partial\rho_j\wedge w\|^2,\\
 \sum_j\|T_s(\rho_jw)\|^2
 &=\|T_sw\|^2+\sum_j\|s\partial\rho_j\wedge w\|^2,\\
 \sum_j\|sP(\rho_jw)\|^2
 &=\|sPw\|^2+\sum_j\|s[P,\rho_j]w\|^2,\\
 \sum_j\int_Y\langle Q(\rho_jw),\rho_jw\rangle_k\,dV_Y
 &=\int_Y\langle Qw,w\rangle_k\,dV_Y.
\end{align*}
For example, the first cross term is
$2\operatorname{Re}\int\langle s\bar\partial w,
s(\sum_j\rho_j\bar\partial\rho_j)\wedge w\rangle=0$;
the other two have the same linear dependence on
$\sum_j\rho_jd\rho_j$.
For each real $\rho_j$, exterior
multiplication and contraction satisfy
\[
 |\bar\partial\rho_j\wedge w|^2
 =|\partial\rho_j|^2|w|^2
 =|\partial\rho_j\wedge w|^2
      +|(\partial\rho_j\wedge)^*w|^2.
\]
Consequently the quadratic partition errors on the two sides cancel,
leaving
\[
 \|\bar\partial w\|_k^2
 \ge\|T_sw\|^2+\|sPw\|^2
             +\int_Y\langle Qw,w\rangle_k\,dV_Y.
\]
The frame covariance just established makes the first two integrals
well defined.  Discarding them proves
\eqref{eq:singular-dual-test}.

Now let $f$ be an $E$-valued $(a,m)$-form of finite inverse-$B$
energy.  It has locally $L^2$ coefficients; closedness is automatic
in this bidegree.  The complementary-index eigenvalue calculation
above gives the pointwise wedge-pairing inequality
\[
 \left|\frac{f\wedge w}{dV_Y}\right|
 \le\langle B^{-1}f,f\rangle_h^{1/2}
             \langle Qw,w\rangle_k^{1/2}.
\]
Indeed, the coefficient indexed by $I$ in $f$ pairs only with that
indexed by $I^c$ in $w$, the two curvature factors are
$\beta_I^{-1}$ and $\beta_I$, and the line-metric factors are
$h$ and $k=h^{-1}$.  Integration and Cauchy--Schwarz give
\[
 \left|\int_Yf\wedge w\right|^2
 \le I_A(f)\int_Y\langle Qw,w\rangle_k\,dV_Y
 \le I_A(f)\|\bar\partial w\|_k^2.
\]
Define a complex-linear functional on the subspace
\[
 \mathcal R=\{\bar\partial w:
  w\in C_c^\infty(Y,\Lambda^{r,0}T_Y^*\otimes E^*)\}
 \subset L^2_{r,1}(Y,E^*,k)
\]
by
\[
 \ell(\bar\partial w)=(-1)^{a+m}\int_Yf\wedge w.
\]
The preceding inequality proves that this definition is independent
of the choice of $w$: a difference with $\bar\partial w=0$ has zero
pairing with $f$.  It also gives $\|\ell\|\le I_A(f)^{1/2}$.
Hahn--Banach extends $\ell$ to the whole Hilbert space with the same
bound.  Riesz representation, followed by the pointwise isometric
wedge duality between $(r,1)$-forms with metric $k$ and
$(a,m-1)$-forms with metric $h=k^{-1}$, produces $u$ satisfying
\[
 \|u\|_h^2\le I_A(f),\qquad
 \int_Yu\wedge\bar\partial w
       =(-1)^{a+m}\int_Yf\wedge w.
\]
The isometry is the usual conjugate Hodge-star identification,
combined with the duality $E\otimes E^*\to\mathbb C$.
Since $h$ is locally bounded below,
$u$ is locally $L^2$ in a smooth metric.  Both sides are therefore
ordinary distributional pairings.  Finally,
\[
 0=\int_Y\bar\partial(u\wedge w)
   =\int_Y\bar\partial u\wedge w
       +(-1)^{a+m-1}\int_Yu\wedge\bar\partial w
\]
is precisely the distributional identity $\bar\partial u=f$.
\end{proof}

\begin{proof}[Proof of Theorem~\ref{thm:singular-product}]
In product coordinate charts, the local weight of $h$ is
$\varphi_G(y)+\varphi_M(z)$.  The multiplier-ideal product formula
\cite[Theorem~2.6(i)]{DEL00} gives
\begin{equation}\label{eq:product-multiplier}
 \mathcal I(h)=
 \bigl(\operatorname{pr}_Y^{-1}\mathcal I(h_G)\mathcal O_X\bigr)
 \bigl(\operatorname{pr}_Z^{-1}\mathcal I(h_M)\mathcal O_X\bigr).
\end{equation}
Put
\[
 \mathcal A=G\otimes\mathcal I(h_G),\qquad
 \mathcal B=\Omega_Z^p\otimes M\otimes\mathcal I(h_M),
\]
Then
\[
 \Omega_{X/Y}^p\otimes L\otimes\mathcal I(h)
       \simeq\mathcal A\boxtimes\mathcal B.
\]
The external tensor product is an ordinary tensor product: the
coherent sheaf $\operatorname{pr}_Z^*\mathcal B$ is flat over $Y$.
Indeed, $\mathcal B$ is flat over the one-point base, and its pullback
to the product is obtained by base change to $Y$; flatness is preserved
under this base change.

Resolve $\mathcal A$ locally on the smooth base by a
bounded complex $P^\bullet$ of finite free modules.  Such a resolution
exists because the local analytic rings of $Y$ are regular.
The complex
\[
 \pi^*P^\bullet\otimes\operatorname{pr}_Z^*\mathcal B
\]
is exact in negative degrees by flatness and resolves
$\mathcal A\boxtimes\mathcal B$.
The derived projection formula is therefore applicable to this
ordinary external product and gives
\[
 R\pi_*(\mathcal A\boxtimes\mathcal B)
 \simeq\mathcal A\otimes^{\mathbf L}_{\mathcal O_Y}
             R\pi_*\operatorname{pr}_Z^*\mathcal B.
\]
Product base change gives the canonical identifications
\[
 R^j\pi_*\operatorname{pr}_Z^*\mathcal B
       \simeq\mathcal O_Y\otimes_{\mathbb C}H^j(Z,\mathcal B).
\]
These sheaves are locally free: their fibers are the fixed
finite-dimensional cohomology groups of the compact manifold $Z$.
In particular, the projection spectral sequence has terms
\[
 E_2^{-i,j}=\operatorname{Tor}^{\mathcal O_Y}_i
 \bigl(\mathcal A,
       \mathcal O_Y\otimes_{\mathbb C}H^j(Z,\mathcal B)\bigr)
 \Longrightarrow
 R^{j-i}\pi_*(\mathcal A\boxtimes\mathcal B),
\]
and $E_2^{-i,j}=0$ for $i>0$.  It collapses to the column $i=0$,
giving
\[
 R^q\pi_*(\mathcal A\boxtimes\mathcal B)
       \simeq\mathcal A\otimes_{\mathbb C}H^q(Z,\mathcal B)
       =\mathcal A\otimes_{\mathbb C}V.
\]
Together with \eqref{eq:product-multiplier}, this proves the first
equality in \eqref{eq:singular-product-sheaf}.

Choose a base point outside the analytic cosupport of
$\mathcal I(h_G)$ at which the weight of $h_G$ is finite.  Such points
have full measure.  On its fiber, the weight differs from the weight
of $h_M$ by the finite constant $\varphi_G(y)$, so
\[
 \mathcal I(h|_{X_y})=\mathcal I(h_M),\qquad
 \|\alpha\|_{h|_{X_y},\omega_Z}^2
       =h_G(y)\|\alpha\|_{h_M,\omega_Z}^2
\]
in a local frame of $G$.  Multiplication of every norm in the fixed
fiber complex by the same positive constant leaves its domains and
its cohomology unchanged.  Taking the minimum over all closed
representatives in Lemma~\ref{prop:fiber-norm} consequently gives
$H_y=h_G(y)H_V$ under the canonical fiber identification.

If $V=0$, the remaining conclusions are immediate.  Otherwise choose
an $H_V$-orthonormal basis $e_1,\ldots,e_N$.  A holomorphic section
$s=\sum_{\nu=1}^N s_\nu\otimes e_\nu$ of $E_0$ has
\[
 |s|_H^2=\sum_{\nu=1}^N|s_\nu|_{h_G}^2.
\]
Since this is a finite sum of nonnegative functions, its local
integrability is equivalent to local integrability of every term.
This proves
$\mathscr E(H)=(G\otimes\mathcal I(h_G))\otimes_{\mathbb C}V$.

Finally write a datum in either stated extremal degree as
$f=\sum_{\nu=1}^N f_\nu\otimes e_\nu$.  The basis is constant and
holomorphic, hence $\bar\partial f=0$ implies
$\bar\partial f_\nu=0$ for every $\nu$.  The auxiliary curvature
acts trivially on $V$, so
\[
 B_A=B_{A,G}\otimes\Id_V,\qquad
 I_A(f)=\sum_{\nu=1}^N I_A(f_\nu)<\infty.
\]
Lemma~\ref{lem:singular-scalar} supplies solutions
$\bar\partial u_\nu=f_\nu$ with
$\|u_\nu\|_{h_Gh_A}^2\le I_A(f_\nu)$.
For $u=\sum_\nu u_\nu\otimes e_\nu$, orthogonality gives
\[
 \bar\partial u=f,\qquad
 \|u\|_{Hh_A}^2
   =\sum_{\nu=1}^N\|u_\nu\|_{h_Gh_A}^2
   \le\sum_{\nu=1}^N I_A(f_\nu)=I_A(f).
\]
This proves both families of estimates for every $p,q$.
Proposition~\ref{prop:scalar-mep}, applied to $H=h_G\otimes H_V$
and $\mathcal F=(G\otimes\mathcal I(h_G))\otimes V$, proves
Griffiths semipositivity and the minimal extension property.
\end{proof}

\begin{remark}\label{rem:singular-scope}
The fiberwise norm identification holds almost everywhere.
At a base point where $h_G=+\infty$, the restricted weight is
identically $-\infty$ on the fiber.  The coherent sheaf in
\eqref{eq:singular-product-sheaf} is defined on all of $Y$.
\end{remark}
\section{Top adjoint descent for arbitrary singular semipositive metrics}
\label{sec:singular-top-adjoint}

Let $\pi:X\to Y$ be a proper surjective holomorphic submersion with
connected fibers of dimension $n\geq1$, where $X$ and $Y$ are connected
complex manifolds. Fix a K\"ahler form $\omega_X$ on $X$ and put
$\omega_y=\omega_X|_{X_y}$. For a singular Hermitian line bundle $(L,h)$
with psh local weights, write
\[
 T_h=\frac{i}{2\pi}\Theta_{L,h},\qquad
 \mathcal F_n=R^n\pi_*
   \bigl(K_{X/Y}\otimes L\otimes\mathcal I(h)\bigr).
\]
Thus $dd^c=(i/2\pi)\partial\bar\partial$ in this section.
This normalization gives $T_{h_D}=[D]$ for the canonical singular
metric $h_D$ on $\mathcal O_X(D)$: if $s_D=f e_D$ is the canonical
section, then $|e_D|_{h_D}^2=|f|^{-2}$ and $|s_D|_{h_D}=1$ off $D$.

The proof of Theorem~\ref{thm:singular-top-rigidity} uses the
following two sheaf-theoretic facts.

\begin{lemma}[Top-degree duality and projection]
\label{lem:top-duality-projection}
For a coherent sheaf $\mathcal A$ on $X$,
\begin{equation}\label{eq:top-duality-coherent}
 \mathcal Hom_Y\!\left(
 R^n\pi_*(K_{X/Y}\otimes\mathcal A),\mathcal O_Y\right)
 \simeq \pi_*\mathcal Hom_X(\mathcal A,\mathcal O_X).
\end{equation}
For a coherent sheaf $\mathcal B$ on $Y$,
\begin{equation}\label{eq:top-projection-coherent}
 R^n\pi_*(K_{X/Y}\otimes\pi^*\mathcal B)
 \simeq\mathcal B.
\end{equation}
\end{lemma}
\begin{proof}
Put $C=R\pi_*(K_{X/Y}\otimes\mathcal A)$, with cohomological indexing,
so that $\mathcal H^k(C)=0$ unless $0\leq k\leq n$; the upper bound
is the coherent cohomological-dimension theorem for a proper map
with $n$-dimensional fibers.  Work locally on the smooth base.
Coherent sheaves on $X$ and $Y$ have finite locally free resolutions
locally.  The smooth morphism $\pi$ has relative dualizing complex
$\pi^!\mathcal O_Y=K_{X/Y}[n]$.  Analytic relative duality
\cite{RRV71} therefore gives
\[
 \begin{aligned}
 R\mathcal Hom_Y(C,\mathcal O_Y)
 &\simeq R\pi_*R\mathcal Hom_X
       (K_{X/Y}\otimes\mathcal A,K_{X/Y}[n])\\
 &\simeq R\pi_*R\mathcal Hom_X(\mathcal A,\mathcal O_X)[n].
 \end{aligned}
\]
The invertible factor $K_{X/Y}$ cancels inside the derived Hom.
Our shift convention is $\mathcal H^j(Q[n])=\mathcal H^{j+n}(Q)$.
On the left the spectral sequence is
\[
 \begin{gathered}
 E_2^{p,q}=\mathcal Ext_Y^p(\mathcal H^{-q}(C),\mathcal O_Y)
 \Longrightarrow
 \mathcal H^{p+q}(R\mathcal Hom_Y(C,\mathcal O_Y)),\\
 p\geq0,\qquad -n\leq q\leq0.
 \end{gathered}
\]
Total degree $-n$ has only $(p,q)=(0,-n)$. An incoming differential
to this term has negative first index, whereas an outgoing differential
$d_r$ has second index $-n-r+1<-n$. Thus the degree $-n$ edge map
identifies its abutment with
$\mathcal Hom_Y(R^n\pi_*(K_{X/Y}\otimes\mathcal A),\mathcal O_Y)$.
On the right, before shifting, the spectral sequence
\[
 {}'E_2^{a,b}=R^a\pi_*\mathcal Ext_X^b(\mathcal A,\mathcal O_X)
 \Longrightarrow
 \mathcal H^{a+b}(R\pi_*R\mathcal Hom_X(\mathcal A,\mathcal O_X))
\]
has $a,b\geq0$. At $(0,0)$ the source of an incoming $d_r$ has
first index $-r$, and the target of an outgoing $d_r$ has second
index $1-r$; both vanish for $r\geq2$.  Its degree-zero edge term is
therefore $\pi_*\mathcal Hom_X(\mathcal A,\mathcal O_X)$.
Taking degree $-n$ after the shift
proves \eqref{eq:top-duality-coherent}.

We record the map underlying this isomorphism.  Normalize the relative
trace
\[
 \operatorname{Tr}_\pi:R^n\pi_*K_{X/Y}\longrightarrow\mathcal O_Y
\]
so that its fiber map is integration of Dolbeault $(n,n)$-classes.  If
$\lambda:\mathcal A|_{\pi^{-1}(U)}\to\mathcal O_{\pi^{-1}(U)}$
is a sheaf morphism, the duality map sends it to
\begin{equation}\label{eq:top-duality-trace-map}
 \operatorname{Tr}_\pi\circ
 R^n\pi_*(\Id_{K_{X/Y}}\otimes\lambda):
 R^n\pi_*(K_{X/Y}\otimes\mathcal A)|_U\longrightarrow\mathcal O_U.
\end{equation}
This is the evaluation pairing followed by the trace in relative
duality.  Its functoriality fixes the cohomological identifications
used below.

The submersion $\pi$ is flat. Since $Y$ is smooth, $\mathcal B$ also has
local finite locally free resolutions, so the derived projection formula gives
\[
 R\pi_*(K_{X/Y}\otimes\pi^*\mathcal B)
 \simeq (R\pi_*K_{X/Y})\otimes^{\mathbf L}\mathcal B.
\]
Here ordinary $\pi^*\mathcal B$ equals its derived pullback. The tensor
spectral sequence has terms
\[
 {}''E_2^{-j,k}
 =\mathcal Tor_j^{\mathcal O_Y}(R^k\pi_*K_{X/Y},\mathcal B),
 \qquad j\geq0,\quad 0\leq k\leq n,
\]
and abuts in degree $k-j$. In degree $n$ its only term is
${}''E_2^{0,n}=(R^n\pi_*K_{X/Y})\otimes\mathcal B$.
An outgoing $d_r$ has positive first index; an incoming $d_r$ starts at
$(-r,n+r-1)$, whose second index exceeds $n$. Both vanish.
Fiberwise Serre duality gives
$h^n(X_y,K_{X_y})=h^0(X_y,\mathcal O_{X_y})=1$ for every $y$.
Grauert's theorem therefore makes $R^n\pi_*K_{X/Y}$ a line bundle
with base change. Its trace map is an isomorphism on every fiber, hence
\[
 R^n\pi_*K_{X/Y}\simeq\mathcal O_Y.
\]
The preceding degree-$n$ edge isomorphism followed by
$\operatorname{Tr}_\pi\otimes\Id_{\mathcal B}$ proves
\eqref{eq:top-projection-coherent}.  In particular, if $\mathcal B$
is a line bundle near $y$, its fiber map is
\[
 H^n(X_y,K_{X_y})\otimes\mathcal B_y\longrightarrow\mathcal B_y,
 \qquad [v]\otimes e\longmapsto
       \left(\int_{X_y}v\right)e.
\]
\end{proof}

\begin{proof}[Proof of Theorem~\ref{thm:singular-top-rigidity}]
On the smooth space $X$, the double dual of the coherent rank-one ideal
$\mathcal I(h)$ is an invertible ideal $\mathcal O_X(-D)$, where $D$ is
effective. Factoring out this invertible ideal gives
\begin{equation}\label{eq:divisorial-ideal-factorization}
 \mathcal I(h)=\mathcal O_X(-D)\mathcal J,
 \qquad
 \operatorname{codim}_X
 \operatorname{Supp}(\mathcal O_X/\mathcal J)\geq2.
\end{equation}
For a prime divisor $P$, its coefficient in $D$ is
$\lfloor\nu_P(T_h)\rfloor$. To verify this for arbitrary psh weights,
put $a=\nu_P(T_h)$. Near a smooth generic point of $P=\{z_1=0\}$,
Siu decomposition gives
\[
 \varphi=a\log|z_1|^2+\eta,\qquad
 \eta\text{ psh},\qquad \nu_P(dd^c\eta)=0.
\]
For $k=\lfloor a\rfloor$, put
\[
 u_k=\tfrac12(\varphi-k\log|z_1|^2).
\]
In the convention $d^c=\tfrac{i}{2\pi}(\bar\partial-\partial)$
used in \cite{Dem12}, the psh function $u_k$ has Lelong number
$a-k<1$ at generic points of $P$.  Skoda's criterion
\cite[Lemma~5.6(a)]{Dem12} therefore gives
\[
 e^{-2u_k}=|z_1|^{2k}e^{-\varphi}\in L^1_{\mathrm{loc}},
 \qquad \operatorname{ord}_P\mathcal I(h)\leq k
\]
there.  The factor $1/2$ converts the weight convention in that
criterion to the convention $h=e^{-\varphi}$ used here.
Conversely, if $g=z_1^\ell u$ with $u$ a unit and $\ell<k$, local
upper bounds for $\eta$ and a lower bound for $|u|$ give
\[
 |g|^2e^{-\varphi}\geq c|z_1|^{-2(a-\ell)},\qquad
 \int_0^\epsilon r^{1-2(a-\ell)}\,dr=+\infty,
 \quad a-\ell\geq1.
\]
Thus $g\notin\mathcal I(h)$ and the generic order is exactly $k$.
Coherence makes $D$ locally finite. Subtracting these integer parts
from the divisorial terms in Siu decomposition gives
\begin{equation}\label{eq:residual-positive-current}
 M=L(-D),\qquad h_M=h/h_D,
 \qquad T_{h_M}=T_h-[D]\geq0.
\end{equation}
The residual weights, initially defined almost everywhere, have psh
representatives because their distributional complex Hessians are
positive.  We use these representatives from now on.
There is also an equality of multiplier ideals
\[
 \mathcal J=\mathcal I(h_M),\qquad
 \mathcal I(h)=\mathcal O_X(-D)\mathcal I(h_M).
\]
Indeed, write $s_D=f e_D$ and
$\varphi_h=\log|f|^2+\varphi_M$ in compatible frames. Then
\[
 |g|^2e^{-\varphi_h}=|g/f|^2e^{-\varphi_M}.
\]
If this is integrable, the local upper bound for the psh function
$\varphi_M$ implies $g/f\in L^2_{\mathrm{loc}}$ off $D$. A meromorphic
function with this property has no divisorial poles: in a transverse
coordinate a pole of order $r\geq1$ has divergent integral
$\int_0^\epsilon t^{1-2r}\,dt$. It then extends across the remaining
codimension-two set by Hartogs. Thus $g/f$ is holomorphic and belongs
to $\mathcal I(\varphi_M)$. The reverse implication follows from the
same displayed equality.

Matsumura's torsion-freeness theorem
\cite[Corollary~1.5]{Mat22} applies because the restrictions of
$\omega_X$ make $\pi$ a proper K\"ahler morphism in the sense of
\cite[Remark~1.4(1)]{Mat22}.  The projection formula and
$K_X=K_{X/Y}\otimes\pi^*K_Y$ give
\[
 \mathcal F_n\simeq
 R^n\pi_*(K_X\otimes L\otimes\mathcal I(h))\otimes K_Y^{-1},
\]
so $\mathcal F_n$ is torsion-free.
Suppose $\mathcal F_n\ne0$. Its generic rank is positive. To obtain a
nonzero dual germ at an arbitrary $y\in Y$, put $R=\mathcal O_{Y,y}$
and $K=\operatorname{Frac}(R)$. Choose a nonzero $K$-linear functional
on $(\mathcal F_n)_y\otimes_R K$. If $v_1,\ldots,v_N$ generate
$(\mathcal F_n)_y$ and $\ell(v_j)=a_j/b_j$, put $d=\prod_jb_j\ne0$.
Then
\[
 0\ne d\ell\in\operatorname{Hom}_R((\mathcal F_n)_y,R).
\]
Coherence extends this germ to a section of $\mathcal F_n^*$ near $y$.
Moreover,
\[
 \mathcal Hom_X(\mathcal J,\mathcal O_X)=\mathcal O_X.
\]
In fact, away from the codimension-two cosupport of $\mathcal J$,
such a homomorphism is multiplication by a holomorphic function.
Hartogs extends that function to $X$, and the extended multiplication
agrees with the homomorphism because $\mathcal O_X$ is torsion-free.
Applying \eqref{eq:top-duality-coherent} to $\mathcal A=M\otimes\mathcal J$
therefore yields
\begin{equation}\label{eq:top-dual-section}
 \mathcal F_n^*\simeq\pi_*M^{-1}.
\end{equation}
After shrinking to a connected neighborhood $U$ of $y$, we obtain
a nonzero $\sigma\in H^0(\pi^{-1}(U),M^{-1})$.
The space $\pi^{-1}(U)$ is connected: the map is proper and open,
and its fibers are connected.  Thus $\sigma$ defines the effective divisor
$E=\operatorname{div}(\sigma)$.

We next show that every component of $E$ is vertical. Cover
$\pi^{-1}(U)$ by countably many relatively compact submersion charts
$V_\alpha\times W_\alpha$, in which $M$ is trivial and the psh weights
$\varphi_\alpha$ are integrable. Fubini gives measure-zero sets
$N_\alpha\subset V_\alpha$ such that
\[
 \begin{gathered}
 \int_{V_\alpha}\int_{W_\alpha}
 |\varphi_\alpha(y,z)|\,dV_z\,dV_y<\infty
 \\
 \Longrightarrow\quad
 \int_{W_\alpha}|\varphi_\alpha(y,z)|\,dV_z<\infty
 \quad(y\notin N_\alpha).
 \end{gathered}
\]
Using a countable exhaustion in each chart, put
$U_0=U\setminus\bigcup_\alpha N_\alpha$. For $y\in U_0$ all restricted
weights are locally integrable psh functions and define a semipositive
metric on $M_y$. Also put
$S=\{y\in U:\dim(E\cap X_y)=n\}$. Properness and the fiber-dimension
theorem make $S$ a proper analytic subset. For $y\in U_0\setminus S$,
the section $\sigma_y$ is nonzero and has effective divisor $E_y$.
Consequently
\begin{equation}\label{eq:zero-fiber-degree}
 0\leq\int_{X_y}T_{h_M|_{X_y}}\wedge\omega_y^{n-1}
 =\int_{X_y}c_1(M_y)\wedge\omega_y^{n-1}
 =-\int_{E_y}\omega_y^{n-1}\leq0.
\end{equation}
For the middle equality, choose a smooth reference metric $h_0$ on
$M_y$ and write $h_M|_{X_y}=h_0e^{-\chi}$, with
$\chi\in L^1(X_y)$. Since $d\omega_y=0$ and $X_y$ is compact,
\[
 \int_{X_y}dd^c\chi\wedge\omega_y^{n-1}
 =\langle\chi,dd^c\omega_y^{n-1}\rangle=0.
\]
The next equality uses $M_y^{-1}\simeq\mathcal O_{X_y}(E_y)$.
If a component of $E$ dominated $U$, its generic fiber would be a
nonzero effective divisor. Positivity of its K\"ahler volume would
contradict \eqref{eq:zero-fiber-degree}. Thus $E$ is vertical.

Let $E_0$ be a vertical prime component and $B_0=\pi(E_0)$. By
properness, $B_0$ is irreducible analytic, and the fiber-dimension bound gives
\[
 \dim B_0\geq\dim E_0-n=\dim U-1.
\]
Hence $B_0$ is a prime divisor. Over its connected smooth locus,
$\pi^{-1}(B_0)$ is a connected complex manifold: the restricted map
is a proper submersion with connected fibers. It is therefore
irreducible. The closed analytic subset $E_0$ has the same dimension
and dominates $B_0$, so it equals the closure of this inverse image.
If $b_0$ locally defines a smooth point of $B_0$, the submersion property
gives $d(b_0\circ\pi)\ne0$ there. Thus the pullback divisor is generically
reduced; no additional divisorial component can lie over the
codimension-two singular locus of $B_0$. Assigning to $B_0$ the
multiplicity of $E_0$ in $E$ proves
\begin{equation}\label{eq:vertical-divisor-pullback}
 E=\pi^*B
\end{equation}
for an effective divisor $B$ on $U$. Properness makes these image
components locally finite: a compact inverse image of a compact
base neighborhood meets only finitely many local components of $E$.
Shrink $U$ so that $B=(b)$.
The meromorphic section $\tau_U=\sigma/(b\circ\pi)$ has order zero
along every prime divisor.  In a local frame, its coefficient and
the reciprocal coefficient are holomorphic off a codimension-two
set and extend by Hartogs.  Thus $\tau_U$ is a nowhere-vanishing
holomorphic frame of $M^{-1}$. In the dual
frame $e_U=\tau_U^{-1}$, write $|e_U|_{h_M}^2=e^{-\varphi_U}$.
The weight $\varphi_U$ is a global psh function on $\pi^{-1}(U)$.

Its restriction to each compact connected fiber is constant or
identically $-\infty$. Define $\psi_U(y)$ to be that value; then
\begin{equation}\label{eq:residual-weight-descends}
 \varphi_U=\psi_U\circ\pi.
\end{equation}
For a local holomorphic section $s:W\to\pi^{-1}(U)$ of the submersion,
$\psi_U|_W=\varphi_U\circ s$ is psh. It cannot be identically
$-\infty$ on an open $W$, since that would make $\varphi_U$ identically
$-\infty$ on the open set $\pi^{-1}(W)$.
On overlaps write $e_V=a_{UV}e_U$. The holomorphic unit $a_{UV}$ is
constant on each compact connected fiber, so
$a_{UV}=g_{UV}\circ\pi$ for a holomorphic unit $g_{UV}$ on $U\cap V$.
This also follows from $\pi_*\mathcal O_X=\mathcal O_Y$.
The frame identities and metric identities give, respectively,
\[
 g_{UV}g_{VW}=g_{UW},\qquad
 \psi_V=\psi_U-\log|g_{UV}|^2.
\]
Thus the $g_{UV}$ define a line bundle $G$, with frames
$\epsilon_V=g_{UV}\epsilon_U$, and
$|\epsilon_U|_{h_G}^2=e^{-\psi_U}$ defines a psh metric on it.
The maps $\pi^*\epsilon_U\mapsto e_U$ glue to
\[
 (M,h_M)\simeq\pi^*(G,h_G),
\]
which proves \eqref{eq:top-singular-factorization}.

The residual multiplier-ideal identity proved above now reads
\[
 \mathcal I(h)=\mathcal O_X(-D)\mathcal I(\psi_U\circ\pi).
\]
In submersion coordinates $(y,z)$ the weight is $\psi_U(y)+0(z)$.
The multiplier-ideal product formula \cite[Theorem~2.6(i)]{DEL00}
therefore gives
\[
 \mathcal I(\psi_U\circ\pi)
 =\pi^{-1}\mathcal I(\psi_U)\mathcal O_X.
\]
These equalities glue, proving the first part of
\eqref{eq:top-singular-sheaf}. In particular there is an isomorphism
of the actual coherent sheaves
\[
 K_{X/Y}\otimes L\otimes\mathcal I(h)
 \simeq K_{X/Y}\otimes\pi^*(G\otimes\mathcal I(h_G)).
\]
Apply \eqref{eq:top-projection-coherent} to
$\mathcal B=G\otimes\mathcal I(h_G)$. This proves
$\mathcal F_n\simeq G\otimes\mathcal I(h_G)$.

We identify this sheaf isomorphism and its fiber maps before computing
the norm.  Put
\[
 \mathcal B=G\otimes\mathcal I(h_G),\qquad
 \mathcal C=K_{X/Y}\otimes L\otimes\mathcal I(h).
\]
Under \eqref{eq:top-singular-factorization}, multiplication by $s_D$
is the holomorphic isomorphism
\begin{equation}\label{eq:top-sheaf-multiplication}
 \mu:K_{X/Y}\otimes\pi^*\mathcal B\xrightarrow{\ \sim\ }\mathcal C,
 \qquad \eta\otimes g\longmapsto s_D\eta\otimes g.
\end{equation}
Locally it sends a coefficient $a$ to $fa$, where $s_D=f e_D$;
the identity $\mathcal I(h)=f\,\pi^*\mathcal I(h_G)$ proves both
surjectivity and injectivity.  Fix the isomorphism
\[
 \Phi:\mathcal F_n\xrightarrow{(R^n\pi_*\mu)^{-1}}
 R^n\pi_*(K_{X/Y}\otimes\pi^*\mathcal B)
 \xrightarrow{\ \operatorname{Tr}_\pi\ }\mathcal B,
\]
where the second arrow is the projection isomorphism of
Lemma~\ref{lem:top-duality-projection}, normalized by its trace.

Define the exceptional sets
\[
 \begin{aligned}
 S_I&=\operatorname{Supp}(\mathcal O_Y/\mathcal I(h_G)),\qquad
 P_G=\{y:\psi_U(y)=-\infty\},\\
 S_D&=\{y:\dim(|D|\cap X_y)=n\}
      =\{y:X_y\subset |D|\},\\
 Y_{\rm good}&=Y\setminus(S_I\cup S_D\cup P_G).
 \end{aligned}
\]
The definition of $P_G$ is independent of the local frame.
The set $S_I$ is proper analytic and $P_G$ has measure zero.
Properness of $D\to Y$ and the fiber-dimension theorem make $S_D$
analytic.  It is proper, since $S_D=Y$ would imply $|D|=X$.
Connectedness of $X_y$ justifies the second expression for $S_D$.
Consequently $Y_{\rm good}$ has full measure.  This choice also
excludes whole fibers contained in horizontal components of $D$.

For $y\in Y_{\rm good}$, the restriction $D_y=D|_{X_y}$ is an
effective Cartier divisor, and $h_G(y)$ is positive and finite.
In compatible local frames, the factorization gives
\[
 \varphi_h=\log|f_j|^2+\psi_U\circ\pi
\]
as an equality of psh functions everywhere: both sides are psh and
agree almost everywhere.  Its restriction to each such fiber gives
\[
 (L_y,h_y)=(\mathcal O_{X_y}(D_y),h_{D_y})\otimes(G_y,h_G(y)).
\]
The preceding divisor integrability argument, applied on $X_y$,
shows that
\begin{equation}\label{eq:top-good-fiber-ideal}
 \mathcal I(h_y)=\mathcal O_{X_y}(-D_y)
               =\mathcal I(h)\mathcal O_{X_y}.
\end{equation}
In the second equality we use $y\notin S_I$ and that a local equation
of $D$ restricts to a nonzero divisor on the smooth connected fiber.
In particular, writing $i_y:X_y\hookrightarrow X$, the natural map
\[
 i_y^*\mathcal C\longrightarrow K_{X_y}\otimes L_y
\]
is locally multiplication by the nonzero divisor $f_j|_{X_y}$.
It is injective, with image
$K_{X_y}\otimes L_y\otimes\mathcal I(h_y)$.
Thus the restriction of \eqref{eq:top-sheaf-multiplication} is
\[
 \mu_y:K_{X_y}\otimes G_y\xrightarrow{\ \sim\ }
 K_{X_y}\otimes L_y\otimes\mathcal I(h_y).
\]

Set $Y'=Y\setminus S_I$.  On $\pi^{-1}(Y')$, the source of $\mu$
is $K_{X/Y}\otimes\pi^*G$.  Naturality of the cohomological
base-change maps gives, for $y\in Y_{\rm good}$, the commutative diagram
\[
 \begin{array}{ccc}
 (R^n\pi_*(K_{X/Y}\otimes\pi^*G))\otimes\mathbb C(y)
 &\xrightarrow{\ (R^n\pi_*\mu)_y\ }&\mathcal F_n\otimes\mathbb C(y)\\[3pt]
 \big\downarrow\beta_y^0&&\big\downarrow\beta_y^h\\[3pt]
 H^n(X_y,K_{X_y})\otimes G_y
 &\xrightarrow{\ H^n(\mu_y)\ }&
 H^n(X_y,K_{X_y}\otimes L_y\otimes\mathcal I(h_y)).
 \end{array}
\]
Both horizontal arrows are isomorphisms.  The map $\beta_y^0$ is
an isomorphism by Grauert's theorem, since
$h^n(X_y,K_{X_y})=1$.  Hence $\beta_y^h$ is an isomorphism as well.
This proves that the fiber space carrying the prescribed norm is
identified with the fiber of the actual direct image by its
canonical base-change map.  Under this identification,
put $\widehat\Phi_y=\Phi_y\circ(\beta_y^h)^{-1}$.  Compatibility
of the relative trace with the fiber trace gives
\begin{equation}\label{eq:top-fiber-trace-identification}
 \widehat\Phi_y\bigl(H^n(\mu_y)([v]\otimes\epsilon)\bigr)
       =\left(\int_{X_y}v\right)\epsilon.
\end{equation}

Fix $y\in Y_{\rm good}$ and $0\ne\epsilon\in G_y$, and write
$c_y=|\epsilon|_{h_G(y)}^2$.  Let
\[
 \mathcal H_r=L^2_{n,r}(X_y,\omega_y),\qquad
 \mathcal K_r=L^2_{n,r}(X_y,L_y;h_y,\omega_y),
 \qquad 0\leq r\leq n.
\]
These spaces consist of equivalence classes modulo sets of measure
zero.  If $s_D=f_j e_j$ and
$u=u_j\otimes(e_j\otimes\epsilon)$ locally, division by $f_j$ on
$X_y\setminus D_y$ defines compatible scalar forms $u_j/f_j$.
Multiplication by $s_D\otimes\epsilon$ is therefore a bijection
of Hilbert spaces
\begin{equation}\label{eq:divisor-isometry-complex}
 T_r:\mathcal H_r\xrightarrow{\ \sim\ }\mathcal K_r,
 \qquad T_rv=s_Dv\otimes\epsilon,\qquad
 \|T_rv\|_{h_y}^2=c_y\|v\|^2.
\end{equation}
Indeed the pointwise equality off $D_y$ is
\[
 |f_jv|^2\,|f_j|^{-2}c_y=c_y|v|^2,
\]
and $D_y$ has measure zero; the inverse map is division by the
local equations followed by any measurable extension on $D_y$.

Let $d_r$ and $d_r^h$ be the maximal distributional
$\bar\partial$ operators on $\mathcal H_r$ and $\mathcal K_r$,
respectively.  Set
$\mathcal H_{n+1}=\mathcal K_{n+1}=0$ and $T_{n+1}=0$.
We verify the equality of graph domains.  For $v\in\operatorname{Dom}d_r$,
the distributional product rule with the holomorphic function $f_j$
gives
\[
 \bar\partial(f_jv)=f_j\bar\partial v,\qquad
 T_rv\in\operatorname{Dom}d_r^h,\qquad
 d_r^hT_rv=T_{r+1}d_rv.
\]
Conversely, take $u\in\operatorname{Dom}d_r^h$ and set
$v=T_r^{-1}u$ and $w=T_{r+1}^{-1}d_r^hu$.  By
\eqref{eq:divisor-isometry-complex}, $v,w$ are ordinary $L^2$ forms,
and off $D_y$,
\[
 \bar\partial v=\bar\partial(u_j/f_j)
               =(\bar\partial u_j)/f_j=w.
\]
Lemma~\ref{lem:l2-removal} gives $\bar\partial v=w$ on the entire
fiber.  Thus
\begin{equation}\label{eq:top-weighted-domain-identities}
 \begin{gathered}
 T_r(\operatorname{Dom}d_r)=\operatorname{Dom}d_r^h,
 \qquad d_r^hT_r=T_{r+1}d_r,\\
 \|T_rv\|_{h_y}^2+\|d_r^hT_rv\|_{h_y}^2
       =c_y\bigl(\|v\|^2+\|d_rv\|^2\bigr).
 \end{gathered}
\end{equation}
In particular $U_r=c_y^{-1/2}T_r$ is unitary, and taking Hilbert
adjoints in $d_r^hU_r=U_{r+1}d_r$ gives, with equality of domains,
\[
 (d_r^h)^*U_{r+1}=U_rd_r^*.
\]
It follows that $T_r$ takes kernels, images, and harmonic subspaces
onto their weighted counterparts.  The ordinary compact Dolbeault
complex has closed range; hence so does the weighted complex.

The induced cohomology map is exactly $H^n(\mu_y)$ in the preceding
base-change diagram.  To see this, use the fine resolutions by
locally $L^2$ maximal graph sheaves in
Lemma~\ref{prop:fiber-norm}.  The same local calculation proves
that $T_\bullet$ is a morphism of these resolutions.  On their
degree-zero kernels its map is
\[
 K_{X_y}\otimes G_y\longrightarrow
 K_{X_y}\otimes L_y\otimes\mathcal I(h_y),\qquad
 \eta\otimes\epsilon\longmapsto s_D\eta\otimes\epsilon,
\]
which is $\mu_y$.  Functoriality of the resolution isomorphism gives
\[
 [T_nv]=H^n(\mu_y)([v]\otimes\epsilon).
\]
Thus the weighted complex realizes the cohomological identification in
\eqref{eq:top-fiber-trace-identification}.

Put $dV_y=\omega_y^n/n!$ and $V_y=\int_{X_y}dV_y$.
The ordinary maximal Dolbeault quotient is $H^n(X_y,K_{X_y})$,
and its trace is $[v]\mapsto\int_{X_y}v$.
For $a\in\operatorname{Dom}d_{n-1}$, distributional Stokes on
the compact fiber gives
\[
 \int_{X_y}d_{n-1}a=0.
\]
Serre duality with $H^0(X_y,\mathcal O_{X_y})=\mathbb C$ makes
this trace an isomorphism.  Thus the representatives of a fixed
class of trace $\lambda$ are precisely the $L^2$ top forms
with integral $\lambda$.
For a unitary coframe $\zeta^1,\ldots,\zeta^n$,
\[
 dV_y=i^{n^2}\zeta^1\wedge\cdots\wedge\zeta^n
           \wedge\overline{\zeta^1}\wedge\cdots
           \wedge\overline{\zeta^n},\qquad |dV_y|^2=1.
\]
Consequently every $v=a\,dV_y$ satisfies
\[
 \left|\int_{X_y}v\right|^2
 =\left|\int_{X_y}a\,dV_y\right|^2
 \leq V_y\int_{X_y}|a|^2\,dV_y=V_y\|v\|^2,
\]
with equality for $v=(\lambda/V_y)dV_y$ when its trace is $\lambda$.
Combining this calculation with the bijection of graph domains
and the cohomological identification yields the exact minimization
formula
\begin{equation}\label{eq:top-quotient-norm-minimization}
 \begin{aligned}
 |\lambda\epsilon|_{H_y}^2
 &=\min_{\widehat\Phi_y([u])=\lambda\epsilon}\|u\|_{h_y}^2\\
 &=c_y\min_{\int_{X_y}v=\lambda}\|v\|^2
   =\frac{c_y|\lambda|^2}{V_y}.
 \end{aligned}
\end{equation}
The minimizing weighted representative is
$s_D(\lambda/V_y)dV_y\otimes\epsilon$.
This expression is independent of the chosen frame: replacing
$\epsilon$ by $a\epsilon$ replaces $c_y$ by $|a|^2c_y$
and $\lambda$ by $a^{-1}\lambda$.

Finally, fiber integration under the proper submersion gives a
smooth function $V_y=\pi_*(\omega_X^n/n!)$, and
\[
 d\,\pi_*(\omega_X^n/n!)=\pi_*d(\omega_X^n/n!)=0.
\]
Connectedness of $Y$ makes it a constant $V>0$.
Equation~\eqref{eq:top-quotient-norm-minimization} proves
$H=V^{-1}h_G$ on $Y_{\rm good}$, and hence almost everywhere,
as asserted in \eqref{eq:top-singular-metric}.
In a frame of $G$, local square integrability for this metric is
exactly $|g|^2e^{-\psi}\in L^1_{\mathrm{loc}}$; the constant $V^{-1}$
does not affect it. Therefore
$\mathscr E(H)=G\otimes\mathcal I(h_G)=\mathcal F_n$.

Finally, for a positive auxiliary line $A$ and a closed $(m,b)$ datum
of finite inverse energy, Lemma~\ref{lem:singular-scalar} gives
$\bar\partial u=f$ and
\[
 \|u\|_{Hh_A}^2=V^{-1}\|u\|_{h_Gh_A}^2
 \leq V^{-1}\int_Y\langle B_A^{-1}f,f\rangle_{h_Gh_A}\,dV_Y
 =\int_Y\langle B_A^{-1}f,f\rangle_{Hh_A}\,dV_Y.
\]
This is the prescribed weighted optimal estimate.
\end{proof}

\begin{remark}\label{rem:top-singular-scope}
For disconnected fibers the argument applies locally to their connected
components, and the resulting metrics form an orthogonal direct sum.
\end{remark}

\section{Counterexamples to unrestricted descent}\label{sec:counterexamples}

\begin{proposition}\label{prop:counter-i}\label{sec:counter-i}
There is a complete K\"ahler product $X=Y\times\mathbb P^1$ and a
smooth metric on $L=\operatorname{pr}_{\mathbb P^1}^*\OO(2)$ such that
\[
 (L,h)\text{ is optimal }(3,2),\qquad
 (\pi_*(K_{X/Y}\otimes L),H)\text{ is not optimal }(2,2).
\]
This is range (i) with $m=2,n=p=1,a=b=2,q=0$.
\end{proposition}
\begin{proof}
Put $r=1/48$, $Y=\Delta_r\times\Delta$, and
\begin{equation}\label{eq:base-example}
 \omega_Y=\frac{i\,dt\wedge d\bar t}{(1-|t|^2/r^2)^2}
          +\frac{i\,ds\wedge d\bar s}{(1-|s|^2)^2},\qquad
 \psi_Y=-r^2\log(1-|t|^2/r^2)-\log(1-|s|^2).
\end{equation}
Differentiation gives
\[
 \partial_t\bar\partial_t[-r^2\log(1-|t|^2/r^2)]
       =(1-|t|^2/r^2)^{-2}.
\]
The analogous identity holds in $s$, so $\omega_Y=i\ddbar\psi_Y$.
Each disk metric is complete: the length of a curve approaching
its boundary dominates the radial integral
$\int_0^R(1-\rho^2/R^2)^{-1}d\rho=+\infty$.
The product metric is therefore complete.  On
$Z=\mathbb P^1$ set
\[
 \omega_Z=\frac{i\,dz\wedge d\bar z}{(1+|z|^2)^2},\quad
 \chi=\frac{1-|z|^2}{1+|z|^2},\quad
 h_0=(1+|z|^2)^{-2}\text{ on }M=K_Z^{-1}.
\]
The expression for $\chi$ in the coordinate $\zeta=1/z$ is
$(|\zeta|^2-1)/(|\zeta|^2+1)$, so it extends smoothly over infinity.
Its coordinate derivatives are
\[
 \chi_z=-\frac{2\bar z}{(1+|z|^2)^2},\qquad
 \chi_{z\bar z}=-\frac{2(1-|z|^2)}{(1+|z|^2)^3}.
\]
Consequently
\begin{equation}\label{eq:chi-identities}
 |\partial\chi|_{\omega_Z}^2
    =\frac{4|z|^2}{(1+|z|^2)^2}\le1,
 \qquad i\ddbar\chi=-2\chi\omega_Z.
\end{equation}
Use $\omega_X=\omega_Y+\omega_Z$ and
\begin{equation}\label{eq:counter-metric}
 h=h_0e^{-6(t+\bar t)\chi}(1-|s|^2)^6.
\end{equation}
Both $X$ and $Y$ admit positive auxiliary lines: use $e^{-\psi_Y}$
on the base and its external product with $\OO_Z(1)$ on $X$.

The weight of $h$ in the standard $\OO(2)$ frame is
\[
 \varphi=2\log(1+|z|^2)+6(t+\bar t)\chi-6\log(1-|s|^2).
\]
Use the unit tangent vectors
\[
 e_t=(1-|t|^2/r^2)\partial_t,\quad
 e_z=(1+|z|^2)\partial_z,\quad
 e_s=(1-|s|^2)\partial_s.
\]
The only mixed coefficient is
\[
 B=\varphi_{t\bar z}(1-|t|^2/r^2)(1+|z|^2)
   =-\frac{12z(1-|t|^2/r^2)}{1+|z|^2},\qquad |B|\le6.
\]
The diagonal coefficients are $0$, $2-12(t+\bar t)\chi$, and $6$.
Thus the curvature matrix is
\begin{equation}\label{eq:counter-matrix}
 \begin{pmatrix}0&B&0\\\bar B&c&0\\0&0&6\end{pmatrix},
 \qquad |B|\le6,\quad c=2-12(t+\bar t)\chi\in(3/2,5/2).
\end{equation}
If $\lambda_-\le\lambda_+$ are the first block's eigenvalues,
$\lambda_-+\lambda_+=c>0$; here
$|12(t+\bar t)\chi|<24r=1/2$ gives the asserted interval for $c$.
Moreover,
\[
 \det\begin{pmatrix}6&B\\\bar B&6+c\end{pmatrix}
       =36+6c-|B|^2\ge6c>0,
\]
The first principal minor of this shifted block is $6>0$, so the
positive determinant makes the block positive definite.  Its
eigenvalues are $6+\lambda_-$ and $6+\lambda_+$, whence
$\lambda_->-6$.  All three pair sums
\begin{equation}\label{eq:positive-pairs}
 c,\qquad6+\lambda_-,\qquad6+\lambda_+
\end{equation}
are positive.  Equations~\eqref{eq:commutator-eigenvalue} and
Proposition~\ref{prop:smooth-test} prove the upstairs estimate.

The holomorphic section $dz\otimes\partial_z$ is the global identity
section of $K_Z\otimes K_Z^{-1}$ and generates its one-dimensional
space of sections.  It has norm one for $(\omega_Z,h_0)$, because
\[
 |dz|_{\omega_Z}^2\,|\partial_z|_{h_0}^2
       =(1+|z|^2)^2(1+|z|^2)^{-2}=1.
\]
In cohomological degree zero the representative is unique, so the
prescribed direct-image metric is its integral:
\[
 H(t,s)=(1-|s|^2)^6\int_Ze^{-6(t+\bar t)\chi}\,\omega_Z.
\]
With $z=\rho e^{i\theta}$,
$\omega_Z=2\rho(1+\rho^2)^{-2}d\rho\,d\theta$ and
$d\chi=-4\rho(1+\rho^2)^{-2}d\rho$.  Therefore the pushforward
of $\omega_Z/(2\pi)$ under $\chi$ is $dx/2$ on $[-1,1]$, giving
\begin{equation}\label{eq:explicit-H}
 H(t,s)=2\pi(1-|s|^2)^6\frac{\sinh u}{u},
 \qquad u=6(t+\bar t).
\end{equation}
Put $F(u)=\log(\sinh u/u)$, with its smooth value at zero.
Then $F''(0)=1/3$ and
\[
 -\partial_t\bar\partial_t\log H(0)=-36F''(0)=-12,\qquad
 -\partial_s\bar\partial_s\log H(0)=6,
 \qquad\partial_t\bar\partial_s\log H=0.
\]
Since $\omega_Y(0)$ is the identity,
\begin{equation}\label{eq:negative-trace}
 i\Theta_H(0)=-12i\,dt\wedge d\bar t+6i\,ds\wedge d\bar s,
 \qquad[i\Theta_H,\Lambda_Y]_{2,2}(0)=-6.
\end{equation}
The necessary implication of Proposition~\ref{prop:smooth-test}
disproves the downstairs estimate.
\end{proof}

\begin{corollary}\label{cor:counter-higher-q}
For every $k\ge1$, a counterexample in range (i) exists with
$m=2,n=k+1,p=n,q=k,a=b=2$, even on a holomorphically trivial family.
\end{corollary}
\begin{proof}
Append factors $Z_j=\mathbb P^1$, $1\le j\le k$, with fixed
Fubini--Study metrics and flat trivial coefficient lines.  Write
$Z_0$ for the original factor, whose adjoint coefficient is trivial.
In the K\"unneth decomposition only the multi-index $(0,1,\ldots,1)$
contributes, since $H^0(\mathbb P^1,K)=0$ and
$H^1(\mathbb P^1,K)=\mathbb C$.  Thus
\[
 H^k\!\left(\prod_{j=0}^kZ_j,K_{\prod Z_j}\otimes L_y\right)
 \simeq H^0(Z_0,K_{Z_0}\otimes L_y)
       \otimes\bigotimes_{j=1}^kH^1(Z_j,K_{Z_j})\simeq\mathbb C.
\]
Let $\eta_j=\omega_{Z_j}/\operatorname{Vol}(Z_j)$.  These are
harmonic representatives of trace one, with
$\|\eta_j\|^2=\operatorname{Vol}(Z_j)^{-1}$.
The fiber Laplacian on the product is the sum of the factor
Laplacians; hence the wedge of the original holomorphic generator
with $\eta_1,\ldots,\eta_k$ is harmonic.  Product integration gives
\[
 H_k(t,s)=H(t,s)\prod_{j=1}^k\operatorname{Vol}(Z_j)^{-1},\qquad
 i\Theta_{H_k}=i\Theta_H.
\]
The curvature eigenvalues upstairs are
$\lambda_-,\lambda_+,6,0,\ldots,0$.  In dimension $N=k+3$, each
$(N,N-1)$ curvature eigenvalue is a sum omitting one of these
numbers.  Omitting $\lambda_-$, $\lambda_+$, $6$, or one of the
zeros gives, respectively,
\[
 6+\lambda_+,\qquad 6+\lambda_-,\qquad c,\qquad6+c,
\]
all positive.  The upstairs estimate follows from
Proposition~\ref{prop:smooth-test}, whereas the downstairs curvature
trace remains $-6$ at the origin.
\end{proof}

The metric \eqref{eq:counter-metric} is not semipositive: its first
curvature block has negative determinant wherever $B\ne0$.

\begin{proposition}\label{prop:counter-ii}\label{sec:counter-ii}
There is a proper holomorphic submersion between complete K\"ahler
manifolds admitting positive auxiliary lines for which
\[
 (\OO_X,1)\text{ is optimal }(1,2),\qquad
 (R^1\pi_*\OO_X,H)\text{ is not optimal }(1,1).
\]
This is range (ii) with $m=n=a=b=q=1,p=0$, and the coefficient
metric is flat.
\end{proposition}
\begin{proof}
Let
\[
 Y=\{|t-t_0|<R\}\Subset\{\operatorname{Im}t>0\},\quad
 \psi_Y=-R^2\log(1-|t-t_0|^2/R^2),\quad\omega_Y=i\ddbar\psi_Y.
\]
Form the holomorphic quotient
\[
 X=(Y\times\mathbb C)/\mathbb Z^2,\qquad
 (a,b)\cdot(t,z)=(t,z+a+bt).
\]
The action is free: $a+bt=0$ and $\operatorname{Im}t>0$ imply
$b=a=0$.  It is properly discontinuous.  Indeed, on any compact
subset of $Y\times\mathbb C$, $\operatorname{Im}t\ge c>0$ and
$|a+bt|\le C$ force
\[
 |b|\le C/c,\qquad |a|\le C+|b|\sup|t|,
\]
so only finitely many lattice translates can meet that compact set.
The quotient is therefore a complex manifold, and the quotient
projection to $Y$ is a holomorphic submersion.
The real coordinates $z=x+tu$, $x,u\in\mathbb R/\mathbb Z$, identify
$X$ smoothly with $Y\times(\mathbb R/\mathbb Z)^2$, proving that
$\pi:X\to Y$ is a proper holomorphic submersion.

Write $v=\operatorname{Im}t$, $w=\operatorname{Im}z$, and $F=w^2/v$.
Its complex Hessian in the order $(t,z)$ is
\[
 (F_{\alpha\bar\beta})
 =\frac1{2v}\begin{pmatrix}w^2/v^2&-w/v\\-w/v&1\end{pmatrix}.
\]
It has rank one and is nonnegative.  Thus, on the cover,
\begin{equation}\label{eq:elliptic-form}
 \Omega=i\ddbar\frac{w^2}{v}
 =\frac{i}{2v}\left(dz-\frac{w}{v}dt\right)
              \wedge\overline{\left(dz-\frac{w}{v}dt\right)}.
\end{equation}
Lattice translation changes the potential by
$2bw+b^2v=\operatorname{Im}(2bz+b^2t)$, so $\Omega$ descends.
It is closed and semipositive, with positive vertical restriction;
therefore $\omega_X=\Omega+\pi^*\omega_Y$ is K\"ahler.  Moreover,
\[
 dz-(w/v)dt=dx+t\,du,\qquad\Omega=dx\wedge du,
 \qquad\int_{X_t}\Omega=1.
\]
For the associated distances,
\[
 \omega_X\ge\pi^*\omega_Y
 \quad\Longrightarrow\quad
 d_Y(\pi(x),\pi(x'))\le d_X(x,x').
\]
If $(x_\nu)$ is $d_X$-Cauchy, then $\pi(x_\nu)\to y\in Y$.
Choose a compact neighborhood $K$ of $y$ containing the projected
tail.  Properness makes $\pi^{-1}(K)$ compact, so some subsequence
converges to $x$.  For any $\varepsilon>0$, choose a sufficiently
late subsequence term with $d_X(x_{\nu_j},x)<\varepsilon/2$ and
use the Cauchy condition to obtain
$d_X(x_\nu,x)\le d_X(x_\nu,x_{\nu_j})+d_X(x_{\nu_j},x)<\varepsilon$.
Thus $(X,\omega_X)$ is complete.

For the positive auxiliary line, take the quotient with automorphy
factor
\[
 (a,b)\cdot(t,z,\xi)
 =(t,z+a+bt,e^{-\pi ib^2t-2\pi ibz}\xi).
\]
If $j_{a,b}(t,z)=e^{-\pi ib^2t-2\pi ibz}$, then
\[
 j_{a,b}(t,z+a'+b't)j_{a',b'}(t,z)
   =e^{-2\pi iba'}j_{a+a',b+b'}(t,z)
   =j_{a+a',b+b'}(t,z).
\]
This is the cocycle identity.  The metric descends because
\[
 |j_{a,b}(t,z)|^2e^{-2\pi(w+bv)^2/v}
  =e^{2\pi b^2v+4\pi bw}e^{-2\pi(w+bv)^2/v}
  =e^{-2\pi w^2/v}.
\]
Its curvature is $i\partial\bar\partial(2\pi w^2/v)=2\pi\Omega$.
Tensoring by the pullback of $e^{-\psi_Y}$ gives positive curvature
$2\pi\Omega+\pi^*\omega_Y$.  For the coefficient line $(\OO_X,1)$,
the curvature commutator is zero, so its upstairs estimate follows
from Lemma~\ref{lem:complete-estimate} after auxiliary twisting.

Under a lattice translation, $dz$ changes to $dz+b\,dt$.
Its class modulo base differentials is therefore invariant, defining
a holomorphic relative form.  On every fiber it spans
$H^0(X_t,K_{X_t})$, so it frames $\pi_*K_{X/Y}$, and
\[
 \omega_t=\frac{i}{2v}dz\wedge d\bar z,\quad
 \int_{X_t}\omega_t=1,\quad\|dz\|^2=2v.
\]
Relative Serre duality gives the holomorphic identification
\[
 R^1\pi_*\OO_X\simeq(\pi_*K_{X/Y})^*.
\]
For the pairing $\int\alpha\wedge dz$, the holomorphic dual frame
has harmonic representative
\[
 \alpha_t=\frac{d\bar z}{2iv},\qquad
 \int_{X_t}d\bar z\wedge dz=2iv,\qquad H(t)=\|\alpha_t\|^2=\frac1{2v}.
\]
The constant-coefficient form $\alpha_t$ is harmonic for the flat
fiber metric.  The equation $\int_{X_t}\alpha_t\wedge dz=1$ shows
that its class is the holomorphic dual frame, even though its
representative depends on $\operatorname{Im}t$.
Explicitly, $|d\bar z|^2=2v$ and the fiber area is one, so
\[
 \|\alpha_t\|^2=(4v^2)^{-1}(2v)\int_{X_t}\omega_t=(2v)^{-1}.
\]
Finally $\partial_t v=-i/2$, $\bar\partial_t v=i/2$, and
$\partial_t\bar\partial_t\log v=-1/(4v^2)$.  Consequently
\[
 i\Theta_H=i\ddbar\log(2v)
              =-\frac{i\,dt\wedge d\bar t}{4v^2}<0.
\]
Proposition~\ref{prop:smooth-test} proves the claimed failure.
\end{proof}


\section{Scalar-factor metrics and minimal extension}\label{sec:extension}

The singular product metric and the metric in the top-adjoint
rigidity theorem have the common form
\begin{equation}\label{eq:scalar-factor-metric}
 E_0=G\otimes_{\mathbb C}V,\qquad
 H=h_G\otimes H_V,\qquad
 \mathcal F=(G\otimes\mathcal I(h_G))\otimes_{\mathbb C}V,
\end{equation}
where $(V,H_V)$ is a fixed finite-dimensional Hermitian space and
$h_G$ has psh weights.  Pointwise values of $H$ are defined using
the psh representatives of these weights.

For a torsion-free sheaf with a singular metric, Griffiths
semipositivity means that $\log|\ell|_{H^*}^2$ has a psh
representative for every local holomorphic section
$\ell\in\mathcal F^*=\mathcal Hom(\mathcal F,\mathcal O_Y)$.
The \emph{minimal extension property} means that there is a proper
analytic $S\subset Y$ such that $\mathcal F$ is locally free off $S$
and, for every open immersion $\iota:B\hookrightarrow Y$ of the
unit ball with $\iota(0)\notin S$ and every $v\in\mathcal F\otimes\mathbb C(\iota(0))$
with $|v|_H=1$, there is $\sigma\in H^0(B,\iota^*\mathcal F)$ with
\[
 \sigma(0)=v,\qquad
 \frac1{\operatorname{Vol}(B)}\int_B|\sigma|_{\iota^*H}^2\,d\lambda\le1.
\]
The measure here is Euclidean; see \cite[Definition 1.2]{Tak24}.

\begin{proposition}\label{prop:scalar-mep}
The metric \eqref{eq:scalar-factor-metric} is Griffiths semipositive
and has the minimal extension property.  Moreover,
$\mathcal F=\mathscr E(H)$, and every holomorphic section on the
complement of a nowhere-dense analytic set with finite $H$-norm on every relatively
compact ambient neighborhood extends uniquely into $\mathcal F$.
\end{proposition}
\begin{proof}
\emph{Integrable sections and removal across analytic sets.}

Fix an $H_V$-orthonormal basis $v_1,\ldots,v_d$ and write
$s=\sum_\nu f_\nu e_G\otimes v_\nu$ in a local frame with
$|e_G|_{h_G}^2=e^{-\varphi}$.  Then
\[
 |s|_H^2=e^{-\varphi}\sum_{\nu=1}^d|f_\nu|^2,
 \qquad
 |s|_H^2\in L^1_{\rm loc}
 \ \Longleftrightarrow\ f_\nu\in\mathcal I(\varphi)
       \quad(1\le\nu\le d).
\]
This proves $\mathscr E(H)=\mathcal F$ coefficientwise.
If $s$ is given off a nowhere-dense analytic subset $A$ and has finite weighted
norm across $A$, the local upper bound $\varphi\le C$ gives
\[
 \int_{K\setminus A}\sum_\nu|f_\nu|^2
 \le e^C\int_{K\setminus A}|s|_H^2<\infty
\]
for compact coordinate neighborhoods $K$.  Extend the coefficients
measurably over $A$, a set of measure zero.  Lemma~\ref{lem:l2-removal}
with $g=0$ gives $\bar\partial f_\nu=0$ distributionally on the
whole neighborhood, hence $f_\nu$ has a holomorphic representative.
For the resulting extension $\widetilde s$,
\[
 \int_K|\widetilde s|_H^2
     =\int_{K\setminus A}|s|_H^2<\infty,
 \qquad \widetilde s\in\mathscr E(H)=\mathcal F.
\]
Holomorphic uniqueness on the dense complement makes these local
extensions glue.

\emph{The dual metric on the actual sheaf.}
To compute the sheaf dual correctly, write locally $h_G=e^{-\varphi}$
and remove the divisorial part of $\mathcal I(\varphi)$:
\[
 \mathcal I(\varphi)=g_D\mathcal J,\qquad
 \psi|_{U\setminus D}=\varphi-\log|g_D|^2,\qquad
 \operatorname{codim}\operatorname{Supp}(\mathcal O/\mathcal J)\ge2.
\]
Here $\psi$ denotes the psh extension across $D$.
These are the divisorial-removal identities in the proof of
Theorem~\ref{thm:singular-top-rigidity}, applied to the local psh
weight on the base.  Put $A=\operatorname{Supp}(\mathcal O/\mathcal J)$.
A morphism $\mathcal J\to\mathcal O$ restricts on $U\setminus A$
to multiplication by a holomorphic function.  Hartogs extension
across $\operatorname{codim}A\ge2$ extends that function to $U$;
equality on $U\setminus A$ then determines the morphism everywhere.
Therefore
\[
 \mathcal Hom(\mathcal J,\mathcal O)=\mathcal O,\qquad
 \mathcal F^{**}=G(-D)\otimes V,\qquad
 \mathcal F^*=G^{-1}(D)\otimes V^*.
\]
In the local frame $g_De_G$ of $G(-D)$, the metric is
$e^{-\psi}H_V$, since $|g_D|^2e^{-\varphi}=e^{-\psi}$.
Consequently, for a dual section $\ell=(\ell_\nu)$,
\[
 \log|\ell|_{H^*}^2
 =\psi+\log\sum_\nu|\ell_\nu|^2,
\]
which is psh.  More explicitly, on the set where
$F=\sum|\ell_\nu|^2>0$, for a tangent vector $\xi$ one has
\[
 (\partial\bar\partial\log F)(\xi,\bar\xi)
 =\frac{F\sum_\nu|\partial_\xi\ell_\nu|^2
          -|\sum_\nu\bar\ell_\nu\partial_\xi\ell_\nu|^2}{F^2}
 \ge0.
\]
The holomorphic logarithmic norm has its psh extension over $F=0$,
and adding the psh weight $\psi$ proves Griffiths semipositivity
on $\mathcal F$.

\emph{Extension with the sharp constant.}
Take $S=\operatorname{Supp}(\mathcal O_Y/\mathcal I(h_G))$.
This is a proper analytic set, and $\mathcal F=E_0$ on $Y\setminus S$.
On an embedded ball $B$, the line bundle $\iota^*G$ is trivial:
the exponential sequence, Cartan's theorem B and contractibility give
\[
 H^1(B,\mathcal O_B^*)\simeq H^2(B,\mathbb Z)=0.
\]
Choose a frame with psh weight $\varphi$.  Write a unit vector at
the center as $v=\sum c_\nu e_G(0)\otimes v_\nu$.  Then
\[
 \varphi(0)>-\infty,\qquad
 e^{-\varphi(0)}\sum_\nu|c_\nu|^2=1.
\]
Sharp scalar extension \cite{GZ15} gives a holomorphic $F$ on $B$
such that
\[
 F(0)=1,\qquad
 \int_B|F|^2e^{-\varphi}\,d\lambda
 \le\operatorname{Vol}(B)e^{-\varphi(0)}.
\]
Define the holomorphic section
$\sigma=\sum c_\nu F e_G\otimes v_\nu$.  It satisfies
\[
 \sigma(0)=v,\qquad
 \int_B|\sigma|_H^2
 =\left(\sum_\nu|c_\nu|^2\right)\int_B|F|^2e^{-\varphi}
 \le\operatorname{Vol}(B).
\]
The same integral bound gives
\[
 c_\nu F\in\mathcal I(\varphi)\quad(1\le\nu\le d),
 \qquad \sigma\in H^0(B,\iota^*\mathcal F).
\]
\end{proof}

\begin{example}[A degreewise estimate does not imply minimal extension]
\label{ex:trace-not-mep}
Take $Y=\mathbb C^2$ with its Euclidean metric,
$Z=\mathbb P^1$, and a trivial fiber coefficient.  On the trivial
base line put
\[
 h_G=e^{-\varphi},\qquad
 \varphi(y)=-|y_1|^2+2|y_2|^2.
\]
For $(m,n,a,b,p,q)=(2,1,2,2,1,1)$ the upstairs curvature
eigenvalues are $-1,2,0$, so $C_{L;3,3}=1$.
The direct image is
$G\otimes H^1(\mathbb P^1,K_{\mathbb P^1})\simeq G$,
with metric $h_G$ up to a constant; it satisfies the optimal
$(2,2)$ estimate, but has a negative curvature direction.

Minimal extension fails at every center, which is necessary to
exclude an exceptional analytic set.  Given $y^0$, pull back to
the unit ball by
$\iota(w)=(y^0_1+\varepsilon w_1,y^0_2+\delta w_2)$,
where $\varepsilon^2>2\delta^2>0$.
The pulled-back weight is
\[
 \varphi\circ\iota=\varphi(y^0)+2\operatorname{Re}\ell(w)
        -\varepsilon^2|w_1|^2+2\delta^2|w_2|^2,
 \qquad
 \ell(w)=-\overline{y^0_1}\varepsilon w_1
                      +2\overline{y^0_2}\delta w_2.
\]
In the holomorphic frame
$e'=e^{\varphi(y^0)/2+\ell(w)}e_G$, the squared norm is
\[
 W(w)=\exp(\varepsilon^2|w_1|^2-2\delta^2|w_2|^2),\qquad W(0)=1.
\]
Any extension of the unit vector $e'(0)$ is $F e'$ with
$F\in\mathcal O(B)$ and $F(0)=1$.  Write
$F=\sum_{\alpha\in\mathbb N^2}c_\alpha w^\alpha$, so $c_0=1$.
For $B_r=\{|w|<r\}$, $0<r<1$, and $\alpha\ne\beta$,
independent coordinate rotations give
\[
 \int_{B_r}w^\alpha\bar w^\beta W\,d\lambda
   =e^{i(\alpha-\beta)\cdot\theta}
           \int_{B_r}w^\alpha\bar w^\beta W\,d\lambda=0
\]
by choosing $\theta$ with $e^{i(\alpha-\beta)\cdot\theta}\ne1$.
Uniform convergence of the Taylor series on $\overline{B_r}$ gives
\[
 \int_{B_r}|F|^2W\,d\lambda
   =\sum_\alpha|c_\alpha|^2\int_{B_r}|w^\alpha|^2W\,d\lambda
   \ge\int_{B_r}W\,d\lambda.
\]
Letting $r\uparrow1$ preserves the inequality by monotone convergence.
Rotational symmetry and polar integration on the unit ball give
$\operatorname{Vol}(B)^{-1}\int_B|w_j|^2\,d\lambda=1/3$.
Jensen's inequality now yields
\[
 \frac1{\operatorname{Vol}(B)}\int_B|F|^2W
 \ge\frac1{\operatorname{Vol}(B)}\int_BW
 \ge\exp\!\left(\frac{\varepsilon^2-2\delta^2}{3}\right)>1.
\]
This contradicts the minimal bound at the arbitrary center $y^0$.
Hence the set of failure is all of $Y$ and cannot be contained in
the exceptional analytic set allowed in the definition.
\end{example}

\begin{thebibliography}{99}
\bibitem[BP08]{BP08}
B.~Berndtsson and M.~P\u aun,
\emph{Bergman kernels and the pseudoeffectivity of relative canonical bundles},
Duke Math. J. \textbf{145} (2008), no.~2, 341--378.
\url{https://doi.org/10.1215/00127094-2008-054}.

\bibitem[BPW22]{BPW22}
Bo Berndtsson, Mihai P\u{a}un, and Xu Wang, \emph{Algebraic fiber spaces and curvature of higher direct images}, J. Inst. Math. Jussieu \textbf{21} (2022), no.~3, 973--1028. \href{https://doi.org/10.1017/S147474802000050X}{doi:10.1017/S147474802000050X}.

\bibitem[Ber09]{Ber09}
B.~Berndtsson,
\emph{Curvature of vector bundles associated to holomorphic fibrations},
Ann. of Math. (2) \textbf{169} (2009), no.~2, 531--560.
\href{https://doi.org/10.4007/annals.2009.169.531}
{doi:10.4007/annals.2009.169.531}.

\bibitem[CP17]{CP17}
J.~Cao and M.~P\u aun,
\emph{Kodaira dimension of algebraic fiber spaces over abelian varieties},
Invent. Math. \textbf{207} (2017), no.~1, 345--387.
\url{https://doi.org/10.1007/s00222-016-0672-6}.

\bibitem[CS24]{CS24}
Y.-J. Choi and G. Schumacher, \emph{Curvature of higher direct images of sheaves of twisted holomorphic forms}, preprint (2024), \href{https://arxiv.org/abs/2407.04298}{arXiv:2407.04298}.

\bibitem[DEL00]{DEL00}
J.-P.~Demailly, L.~Ein, and R.~Lazarsfeld,
\emph{A subadditivity property of multiplier ideals},
Michigan Math. J. \textbf{48} (2000), 137--156.
\href{https://arxiv.org/abs/math/0002035}{arXiv:math/0002035v2};
see Theorem 2.6(i) for the analytic product formula.

\bibitem[DNWZ23]{DNWZ23}
Fusheng Deng, Jiafu Ning, Zhiwei Wang, and Xiangyu Zhou, \emph{Positivity of holomorphic vector bundles in terms of $L^p$-estimates for $\bar\partial$}, Math. Ann. \textbf{385} (2023), nos.~1--2, 575--607. \href{https://doi.org/10.1007/s00208-021-02348-7}{doi:10.1007/s00208-021-02348-7}.

\bibitem[DPS01]{DPS01}
J.-P.~Demailly, T.~Peternell, and M.~Schneider,
\emph{Pseudo-effective line bundles on compact K\"ahler manifolds},
Internat. J. Math. \textbf{12} (2001), no.~6, 689--741.
\url{https://doi.org/10.1142/S0129167X01000861}.

\bibitem[Dem12]{Dem12}
J.-P.~Demailly,
\emph{Analytic Methods in Algebraic Geometry},
Surveys of Modern Mathematics, vol.~1,
International Press and Higher Education Press, 2012.

\bibitem[Dem82]{Dem82}
J.-P.~Demailly,
\emph{Estimations \(L^2\) pour l'op\'erateur \(\bar\partial\) d'un
fibr\'e vectoriel holomorphe semi-positif au-dessus d'une vari\'et\'e
k\"ahl\'erienne compl\`ete},
Ann. Sci. \'Ecole Norm. Sup. (4) \textbf{15} (1982), no.~3, 457--511.
\href{https://www.numdam.org/article/ASENS_1982_4_15_3_457_0.pdf}
{Full text}; see \S5, Theorem 5.1.

\bibitem[Fuj78]{Fuj78}
T. Fujita, \emph{On K\"ahler fiber spaces over curves}, J. Math. Soc. Japan \textbf{30} (1978), no.~4, 779--794. \href{https://doi.org/10.2969/jmsj/03040779}{doi:10.2969/jmsj/03040779}.

\bibitem[GZ15]{GZ15}
Q.~Guan and X.~Zhou,
\emph{A solution of an \(L^2\) extension problem with an optimal
estimate and applications},
Ann. of Math. (2) \textbf{181} (2015), no.~3, 1139--1208.
\href{https://doi.org/10.4007/annals.2015.181.3.6}
{doi:10.4007/annals.2015.181.3.6}.

\bibitem[Gra60]{Gra60}
Hans Grauert, \emph{Ein Theorem der analytischen Garbentheorie und die Modulr\"aume komplexer Strukturen}, Inst. Hautes \'{E}tudes Sci. Publ. Math. \textbf{5} (1960), 5--64. \href{https://doi.org/10.1007/BF02684746}{doi:10.1007/BF02684746}.

\bibitem[HPS18]{HPS18}
C.~Hacon, M.~Popa, and C.~Schnell,
\emph{Algebraic fiber spaces over abelian varieties: around a recent
theorem by Cao and P\u aun},
in \emph{Local and Global Methods in Algebraic Geometry},
Contemp. Math. \textbf{712}, Amer. Math. Soc., Providence, RI, 2018,
143--195.
\url{https://doi.org/10.1090/conm/712/14346}.

\bibitem[IMW24]{IMW24}
T.~Inayama, S.-i.~Matsumura, and Y.~Watanabe,
\emph{Singular Nakano positivity of direct image sheaves of adjoint bundles},
Anal. PDE, to appear;
\url{https://arxiv.org/abs/2407.11412}, version~2 (2024).

\bibitem[Ina26]{Ina26}
T.~Inayama,
\emph{Partial positivity is not preserved by marginalization},
preprint (2026),
\url{https://arxiv.org/abs/2609.15075}.

\bibitem[Kaw81]{Kaw81}
Y. Kawamata, \emph{Characterization of abelian varieties}, Compositio Math. \textbf{43} (1981), no.~2, 253--276.

\bibitem[Kol86a]{Kol86a}
J. Koll\'ar, \emph{Higher direct images of dualizing sheaves. I}, Ann. of Math. (2) \textbf{123} (1986), no.~1, 11--42. \href{https://doi.org/10.2307/1971351}{doi:10.2307/1971351}.

\bibitem[Kol86b]{Kol86b}
J. Koll\'ar, \emph{Higher direct images of dualizing sheaves. II}, Ann. of Math. (2) \textbf{124} (1986), no.~1, 171--202.

\bibitem[LY14]{LY14}
Kefeng Liu and Xiaokui Yang, \emph{Curvatures of direct image sheaves of vector bundles and applications}, J. Differential Geom. \textbf{98} (2014), no.~1, 117--145. \href{https://arxiv.org/abs/1303.2701}{arXiv:1303.2701}.

\bibitem[MT07]{MT07}
C.~Mourougane and S.~Takayama,
\emph{Hodge metrics and positivity of direct images},
J. Reine Angew. Math. \textbf{606} (2007), 167--178.
\url{https://doi.org/10.1515/CRELLE.2007.039}.

\bibitem[MT08]{MT08}
Ch.~Mourougane and S.~Takayama,
\emph{Hodge metrics and the curvature of higher direct images},
Ann. Sci. \'Ecole Norm. Sup. (4) \textbf{41} (2008), no.~6, 905--924.
\href{https://arxiv.org/abs/0707.3551}{arXiv:0707.3551v2};
\href{https://doi.org/10.24033/asens.2084}{doi:10.24033/asens.2084}.

\bibitem[MT09]{MT09}
C.~Mourougane and S.~Takayama,
\emph{Extension of twisted Hodge metrics for K\"ahler morphisms},
J. Differential Geom. \textbf{83} (2009), no.~1, 131--161.
\url{https://arxiv.org/abs/0809.3221}.

\bibitem[Mat22]{Mat22}
S.-i.~Matsumura,
\emph{Injectivity theorems with multiplier ideal sheaves for higher
direct images under K\"ahler morphisms},
Algebr. Geom. \textbf{9} (2022), no.~2, 122--158.
\url{https://doi.org/10.14231/AG-2022-005}.

\bibitem[Nau21]{Nau21}
Philipp Naumann, \emph{Curvature of higher direct images}, Ann. Fac. Sci. Toulouse Math. (6) \textbf{30} (2021), no.~1, 171--201. \href{https://doi.org/10.5802/afst.1670}{doi:10.5802/afst.1670}.

\bibitem[PT18]{PT18}
M.~P\u aun and S.~Takayama,
\emph{Positivity of twisted relative pluricanonical bundles and their
direct images},
J. Algebraic Geom. \textbf{27} (2018), no.~2, 211--272.
\url{https://doi.org/10.1090/jag/702}.

\bibitem[RRV71]{RRV71}
J.-P.~Ramis, G.~Ruget, and J.-L.~Verdier,
\emph{Dualit\'e relative en g\'eom\'etrie analytique complexe},
Invent. Math. \textbf{13} (1971), 261--283.
\url{https://doi.org/10.1007/BF01406078}.

\bibitem[Siu74]{Siu74}
Y.-T.~Siu,
\emph{Analyticity of sets associated to Lelong numbers and the extension
of closed positive currents},
Invent. Math. \textbf{27} (1974), 53--156.
\url{https://doi.org/10.1007/BF01389965}.

\bibitem[Tak24]{Tak24}
S. Takayama, \emph{Singular Griffiths semi-positivity of higher direct images}, Math. Ann. \textbf{388} (2024), 4343--4353. \href{https://doi.org/10.1007/s00208-023-02632-8}{doi:10.1007/s00208-023-02632-8}.

\bibitem[Take95]{Take95}
K.~Takegoshi,
\emph{Higher direct images of canonical sheaves tensorized with
semi-positive vector bundles by proper K\"ahler morphisms},
Math. Ann. \textbf{303} (1995), 389--416.

\bibitem[Vie83]{Vie83}
E. Viehweg, \emph{Weak positivity and the additivity of the Kodaira dimension for certain fibre spaces}, in \emph{Algebraic varieties and analytic varieties (Tokyo, 1981)}, Adv. Stud. Pure Math., vol.~1, North-Holland, Amsterdam, 1983, pp.~329--353. \href{https://doi.org/10.2969/aspm/00110329}{doi:10.2969/aspm/00110329}.

\bibitem[Wat25a]{Wat25}
Y. Watanabe, \emph{Dual Nakano positivity and singular Nakano positivity of direct image sheaves}, Nagoya Math. J. \textbf{257} (2025), 55--78. \href{https://doi.org/10.1017/nmj.2024.20}{doi:10.1017/nmj.2024.20}.

\bibitem[Wat25b]{WatCurv}
Yuta Watanabe, \emph{Curvature operator of holomorphic vector bundles and $L^2$-estimate condition for $(n,q)$ or $(p,n)$-forms}, Tohoku Math. J. (2) \textbf{77} (2025), no.~4, 445--471. \href{https://doi.org/10.2748/tmj.20231128}{doi:10.2748/tmj.20231128}. Theorem numbers cited here refer to \href{https://arxiv.org/abs/2109.12554v2}{arXiv:2109.12554v2}.

\bibitem[Wat25c]{WatNadel}
Y.~Watanabe,
\emph{Nakano--Nadel type, Bogomolov--Sommese type vanishing and
singular dual Nakano semi-positivity},
Ann. Fac. Sci. Toulouse Math. (6) \textbf{34} (2025), no.~2, 339--394.
\url{https://doi.org/10.5802/afst.1815};
\href{https://arxiv.org/abs/2209.00823v2}{arXiv:2209.00823v2}.
The weighted Dolbeault resolution is Theorem 5.3(a) in the cited
arXiv version.

\bibitem[Wu21]{Wu21}
X.~Wu,
\emph{On the hard Lefschetz theorem for pseudoeffective line bundles},
Internat. J. Math. \textbf{32} (2021), no.~6, Paper No.~2150035, 25~pp.
\url{https://doi.org/10.1142/S0129167X2150035X}.
\end{thebibliography}
\end{document}
